\documentclass[aspectratio=169,10pt]{beamer}
% \NOTESONLY (defined before this file is read) prints only the speaker notes
\ifdefined\NOTESONLY
  \setbeameroption{show only notes}
\fi

\usepackage{lmodern}
\usepackage[utf8]{inputenc}
\usepackage[T1]{fontenc}
\usepackage[polish,english]{babel}
\usepackage{amsmath,amssymb}
\usepackage{booktabs}
\usepackage{colortbl}
\usepackage{xcolor}
\usepackage{fontawesome5}
\usepackage{etoolbox}
\usepackage{xstring}
\usepackage{tikz}
\usetikzlibrary{arrows.meta, calc, positioning, fit, decorations.pathreplacing, decorations.pathmorphing}
\usepackage{pgfplots}
% layers for the code windows drawn in TikZ: the window, then the line that runs now, then the text
\pgfdeclarelayer{codebg}\pgfdeclarelayer{codehl}\pgfsetlayers{codebg,codehl,main}
\pgfplotsset{compat=1.18}

\usetheme{Madrid}
\useinnertheme{rounded}
\useoutertheme{miniframes}
\setbeamertemplate{navigation symbols}{}
\setbeamertemplate{itemize items}[circle]
\setbeamertemplate{enumerate items}[default]
\setbeamertemplate{blocks}[rounded][shadow=true]
% one dot per idea, not per frame: \samedot before a frame gives it the dot of the frame before
% (it undoes the step of the dot counter that the next frame makes)
\newcommand{\samedot}{\addtocounter{subsectionslide}{-1}}
% "RC4 fix N/3": an orange label at the right end of the title bar, on the slides that fix RC4 problem N
% (the problems are numbered in Section 1 in the same order: 1 the key is the only input, 2 not random, 3 no integrity)
\newcommand{\fixtag}[1]{%
  \begin{tikzpicture}[remember picture, overlay]
    \node[fill=AccentOrange, text=white, font=\scriptsize\bfseries, rounded corners=3pt, inner sep=3pt, anchor=east]
      at ([xshift=-0.35cm, yshift=-1.1cm]current page.north east) {RC4 fix #1/3};
  \end{tikzpicture}}
% the default dot step of the circle dots (0.1 cm dots)
\setbeamersize{mini frame size=.14cm,mini frame offset=-.01cm}

% ======================================================================
% Colour rules for presentations (Beamer + TikZ), from Agents/colors
% ======================================================================

% ---------- template ----------
\definecolor{DeepBlue}{HTML}{283593}      % title bar, headings, links, control-flow arrows
\definecolor{AccentOrange}{HTML}{EB5A14}  % annotations: step numbers, "look here", code highlight, conclusions
\definecolor{LightGray}{HTML}{F5F6FA}     % background of information blocks
\setbeamercolor{palette primary}{bg=DeepBlue,fg=white}
\setbeamercolor{palette secondary}{bg=DeepBlue!80!black,fg=white}
\setbeamercolor{palette tertiary}{bg=DeepBlue!90!black,fg=white}
\setbeamercolor{title}{bg=DeepBlue,fg=white}
\setbeamercolor{frametitle}{bg=DeepBlue,fg=white}
\setbeamercolor{structure}{fg=DeepBlue}
\setbeamercolor{item}{fg=AccentOrange}
\setbeamercolor{alerted text}{fg=AccentOrange}
\setbeamercolor{block title}{bg=DeepBlue,fg=white}
\setbeamercolor{block body}{bg=LightGray,fg=black}
\setbeamercolor{block title example}{bg=AccentOrange,fg=white}
\setbeamercolor{block body example}{bg=AccentOrange!10,fg=black}
\setbeamercolor{enumerate item}{fg=DeepBlue}
\setbeamerfont{enumerate item}{series=\bfseries}

% ---------- actors and categories: fills only (S1 to S6) ----------
\definecolor{sIndigo}{HTML}{5C6BC0}   % S1
\definecolor{sPink}{HTML}{EC407A}     % S2
\definecolor{sViolet}{HTML}{7E57C2}   % S3
\definecolor{sAmber}{HTML}{FFB300}    % S4: black text on it
\definecolor{sTeal}{HTML}{26A69A}     % S5
\definecolor{sBrown}{HTML}{8D6E63}    % S6
% text variants that stay readable on white
\definecolor{sPinkText}{HTML}{C2185B}
\definecolor{sAmberText}{HTML}{A06A00}
\definecolor{sTealText}{HTML}{00796B}

% ---------- actions: frames only; results: signs only ----------
\definecolor{aExec}{HTML}{43A047}     % runs now
\definecolor{aRead}{HTML}{1E88E5}     % is read
\definecolor{aWrite}{HTML}{E53935}    % is written

% ---------- neutral (Blue Grey) ----------
\definecolor{nInk}{HTML}{37474F}      % text in figures, waveforms, wires, data-flow arrows
\definecolor{nLine}{HTML}{90A4AE}     % borders, axes
\definecolor{nLight}{HTML}{CFD8DC}    % inactive lines, grids, empty places (dashed)
\definecolor{nFill}{HTML}{ECEFF1}     % fill of boxes and used cells
\definecolor{nIdle}{HTML}{E0E0E0}     % time lost (idle, busy waiting, bubbles)

% ---------- TikZ styles ----------
\tikzset{
  aexec/.style={draw=aExec, line width=1.6pt},
  aexecm/.style={aexec, label={[text=aExec, inner sep=1pt, font=\tiny]left:$\blacktriangleright$}},
  aread/.style={draw=aRead, line width=1.6pt, dash pattern=on 3pt off 2pt},
  awrite/.style={draw=aWrite, line width=2.4pt},
  hw/.style={draw=nLine, thick, fill=nFill, text=nInk, rounded corners=3pt, align=center, inner sep=4pt},
  hwmain/.style={draw=nInk, thick, fill=nInk, text=white, rounded corners=3pt, align=center, inner sep=4pt},
  cell/.style={draw=nLine, thick, fill=white, text=nInk, font=\ttfamily\scriptsize},
  cellused/.style={cell, fill=nFill},
  cellfree/.style={cell, dashed, draw=nLight},
  container/.style={draw=nLine, thick, dashed, rounded corners=5pt},
  actor/.style={fill=#1, draw=#1, text=white, rounded corners=3pt, align=center, font=\scriptsize\bfseries, inner sep=4pt},
  actorlight/.style={fill=#1!15, draw=#1, thick, text=nInk, rounded corners=3pt, align=center, inner sep=4pt},
  suspended/.style={fill=#1!30, draw=#1, dashed, line width=0.9pt, text=nInk, rounded corners=3pt, align=center, inner sep=4pt},
  ready/.style={fill=white, draw=#1, line width=0.8pt, text=nInk, rounded corners=3pt, align=center, inner sep=4pt},
  idle/.style={fill=nIdle, draw=none, text=nInk, rounded corners=3pt, align=center, inner sep=4pt},
  flow/.style={-{Stealth}, very thick, DeepBlue},
  data/.style={-{Stealth}, very thick, nInk},
  holds/.style={-{Stealth}, very thick, nInk},
  waits/.style={-{Stealth}, very thick, nLine, dashed},
  annot/.style={draw=AccentOrange, line width=1.4pt, rounded corners=3pt},
}
\tikzset{actoramber/.style={actor=sAmber, text=black}}

% ---------- signs and small helpers ----------
\newcommand{\ok}{\textcolor{aExec}{\faCheck}}
\newcommand{\nok}{\textcolor{aWrite}{\faTimes}}
\newcommand{\badge}[1]{\tikz[baseline=(b.base)]\node[circle,fill=AccentOrange,text=white,
  inner sep=1.2pt,font=\bfseries\tiny] (b) {#1};}
\newcommand{\bigbadge}[1]{\tikz[baseline=(b.base)]\node[circle,fill=AccentOrange,text=white,
  inner sep=2pt,font=\bfseries\small] (b) {#1};}
% ---------- text boxes ----------
\newcommand{\defbox}[1]{%
  \begin{beamercolorbox}[rounded=true,sep=0.1cm,wd=\linewidth]{block body example}
    \footnotesize #1
  \end{beamercolorbox}\vspace{0.06cm}}
% a large definition in the middle of the slide
\newcommand{\bigdef}[1]{%
  \begin{beamercolorbox}[rounded=true,sep=0.25cm,wd=\linewidth]{block body example}
    \centering\normalsize #1
  \end{beamercolorbox}\vspace{0.06cm}}
\newcommand{\stepbox}[2]{%
  \begin{beamercolorbox}[rounded=true,sep=0.09cm,wd=\linewidth]{block body}
    \footnotesize\textbf{#1}\\ #2
  \end{beamercolorbox}\vspace{0.06cm}}
\newcommand{\keymsg}[1]{%
  \begin{beamercolorbox}[rounded=true,sep=0.1cm,wd=\textwidth]{block body example}
    \centering\footnotesize #1
  \end{beamercolorbox}}

% ---------- code: dark window, grey dots, syntax colours ----------
\definecolor{VSDarkBg}{RGB}{25,25,25}
\definecolor{VSText}{RGB}{245,245,245}
\definecolor{VSBlue}{RGB}{80,180,255}
\definecolor{VSPurple}{RGB}{210,110,210}
\definecolor{VSString}{RGB}{255,140,80}
\definecolor{VSComment}{RGB}{80,220,80}
\setbeamercolor{vscodelisting}{bg=VSDarkBg}

% ======================================================================
% Course macros
% ======================================================================
% the starter files on GitHub: the word "GitHub" is the link.
%   \ghlink{Task4} opens the folder of Task 4, \ghlink{} opens the folder of this lecture
\newcommand{\ghbase}{https://github.com/KsaweryM/Cryptography-Course/tree/main/Lecture02-ChaCha20}
\newcommand{\ghlink}[1]{\href{\ghbase/#1}{\textcolor{DeepBlue}{\textbf{GitHub}}}}
% the task slides: the steps are in the README on GitHub, the slide only opens the task
%   \taskgoal{one line}   \taskready{the ready modules}   \opentask{Task2}{about 6 minutes}
\newcommand{\taskgoal}[1]{{\centering\large\textcolor{AccentOrange}{\faBullseye}\ \textbf{Goal:} #1\par}\vspace{0.2cm}}
\newcommand{\taskready}[1]{{\centering\normalsize\textcolor{nInk}{\faBox}\ \textbf{Ready:} #1\par}\vspace{0.2cm}}
\newcommand{\opentask}[2]{%
  \begin{beamercolorbox}[rounded=true,sep=0.35cm,wd=\linewidth]{block body example}
    \centering
    {\LARGE\faGithub\ \textbf{Open the task:}\ \href{\ghbase/#1}{\textcolor{DeepBlue}{\textbf{\texttt{#1}}}}}\\[8pt]
    \footnotesize\href{\ghbase/#1}{\textcolor{DeepBlue}{\texttt{github.com/KsaweryM/Cryptography-Course/tree/main/Lecture02-ChaCha20/#1}}}\\[6pt]
    \small on \ghlink{#1}: the steps, the starter files and how to run them\ifx&#2&\else\qquad\faClock[regular]\ #2\fi
  \end{beamercolorbox}}
% a page of the simulator website, printed in full so that students can type it:
%   \simlink{chacha20/\#2} prints and links ksawerym.github.io/Cryptography-simulators/chacha20/#2
\newcommand{\simbase}{ksawerym.github.io/Cryptography-simulators}
\newcommand{\simlink}[1]{\href{https://\simbase/#1}{\textcolor{DeepBlue}{\texttt{\simbase/#1}}}}
% a slide that sends the room to the simulator (shown live by the lecturer): \opensim{chacha20/\#5}
\newcommand{\opensim}[1]{%
  \begin{beamercolorbox}[rounded=true,sep=0.35cm,wd=\linewidth]{block body example}
    \centering
    {\LARGE\faDesktop\ \textbf{Open the simulator}}\\[8pt]
    \normalsize\simlink{#1}
  \end{beamercolorbox}}
% "For the curious" line under a frame: \curious{chacha20/\#2}
\newcommand{\curious}[1]{{\footnotesize\textcolor{gray}{\faSearch\ For the curious:}\ \simlink{#1}}}

% speaker notes (Polish): printed only in the notes PDF (make notes)
\setbeamertemplate{note page}[default]
\setbeamerfont{note page}{size=\normalsize}
\newcommand{\say}[1]{\note{\selectlanguage{polish}#1}}

% ======================================================================
% Colour mapping for this lecture (colour rules, section 6)
%   S1 sIndigo : the sender (the sensor)
%   S2 sPink   : the attacker
%   S3 sViolet : plaintext
%   S4 sAmber  : key and keystream
%   S5 sTeal   : the receiver (the server)
%   S6 sBrown  : ciphertext
%   mechanisms (generator, block cipher, XOR, RC4 itself) are neutral
% ======================================================================
\tikzset{
  byte/.style={draw=nLine, thick, fill=white, text=nInk, font=\ttfamily\scriptsize,
    minimum width=0.62cm, minimum height=0.48cm, inner sep=1pt, rounded corners=1.5pt},
  ptb/.style={byte, fill=sViolet!15, draw=sViolet},
  ksb/.style={byte, fill=sAmber!25, draw=sAmber},
  ctb/.style={byte, fill=sBrown!15, draw=sBrown},
  % the same ciphertext value always gets the same shade of S6 (lightness, not a new colour)
  ct58/.style={ctb, fill=sBrown!12},
  ct137/.style={ctb, fill=sBrown!35},
  ct201/.style={ctb, fill=sBrown!70, text=white},
  ct64/.style={ctb, fill=sBrown!100, text=white},
  rowlabel/.style={font=\scriptsize, text=nInk, anchor=east},
  sender/.style={actor=sIndigo, minimum width=2.4cm, minimum height=1.0cm, font=\small\bfseries},
  receiver/.style={actor=sTeal, minimum width=2.4cm, minimum height=1.0cm, font=\small\bfseries},
  attacker/.style={actor=sPink, minimum width=2.0cm, minimum height=0.8cm, font=\small\bfseries},
  radio/.style={-{Stealth}, very thick, nInk, decorate,
    decoration={snake, amplitude=1.2pt, segment length=7pt, post length=4pt}},
}
% a row of bytes: \byterow{x of the first cell}{y}{style}{list of values}{step between cells}
\newcommand{\byterow}[5]{%
  \foreach \v [count=\k from 0] in {#4} {
    \node[#3] at ({#1+#5*\k},#2) {\v};
  }}


% ======================================================================
% Helpers shared with Lecture 1 (same look in every lecture)
% ======================================================================
% an intercepted envelope with ciphertext: \envelope{name}{position}{label}
\newcommand{\envelope}[3]{%
  \node[draw=sBrown, thick, fill=sBrown!15, minimum width=2.2cm, minimum height=1.3cm, rounded corners=1pt] (#1) at (#2) {};
  \draw[sBrown, thick] (#1.north west) -- ([yshift=-0.45cm]#1.north) -- (#1.north east);
  \node[circle, fill=sBrown, text=white, inner sep=1.5pt, font=\scriptsize] at ([yshift=-0.45cm]#1.north) {\faLock};
  \node[font=\scriptsize\bfseries, text=nInk] at ([yshift=-0.28cm]#1.center) {#3};}
% code drawn in TikZ, so that single lines can be marked:
%   \codebox{number of lines}{width}  then  \cl{line}{text}, \clmark{line} (runs now), \clnote{line}{text}
\newcommand{\codepitch}{0.36}
\newcommand{\codebox}[2]{%
  \begin{pgfonlayer}{codebg}
    \fill[VSDarkBg, rounded corners=4pt] (-0.15,0.55) rectangle ({#2+0.45},{-\codepitch*#1-0.15});
    % the gutter with the line numbers, as in an editor
    \fill[white, opacity=0.05] (-0.15,0.12) rectangle (0.4,{-\codepitch*#1-0.06});
  \end{pgfonlayer}
  \foreach \d in {0,0.3,0.6} \fill[nLine] ({0.12+\d},0.33) circle (2.4pt);
  \coordinate (codeR) at ({#2+0.45},0);}
\newcommand{\cl}[3][VSText]{%
  \node[anchor=east, font=\ttfamily\tiny, text=VSText!45!VSDarkBg, inner sep=1pt] at (0.34,{-\codepitch*#2+0.14}) {#2};
  \node[anchor=west, font=\ttfamily\scriptsize, text=#1, inner sep=1pt] (L#2) at (0.5,{-\codepitch*#2+0.14}) {\strut #3};}
% the line that runs now, as in a debugger stopped at a breakpoint: the whole line lit from edge to edge,
% an accent on the left edge and an arrow in the gutter. The same mark on every slide that steps through code.
\newcommand{\clmark}[1]{%
  \begin{pgfonlayer}{codehl}
    \fill[AccentOrange!30!VSDarkBg] (-0.15,{-\codepitch*#1+0.14-0.18}) rectangle (codeR |- 0,{-\codepitch*#1+0.14+0.18});
    \fill[AccentOrange] (-0.15,{-\codepitch*#1+0.14-0.18}) rectangle (-0.08,{-\codepitch*#1+0.14+0.18});
  \end{pgfonlayer}
  \node[text=AccentOrange, font=\scriptsize, inner sep=0pt] at (0.04,{-\codepitch*#1+0.14}) {$\blacktriangleright$};}
% a line we talk about (not running): lit a little, from edge to edge, no frame
\newcommand{\clhl}[1]{%
  \begin{pgfonlayer}{codehl}
    \fill[white!14!VSDarkBg] (-0.15,{-\codepitch*#1+0.14-0.18}) rectangle (codeR |- 0,{-\codepitch*#1+0.14+0.18});
  \end{pgfonlayer}}
\newcommand{\clnote}[2]{%
  \clhl{#1}
  \draw[->, thick, AccentOrange, >=Stealth] ([xshift=0.6cm]codeR |- 0,{-\codepitch*#1+0.14}) -- ([xshift=0.08cm]codeR |- 0,{-\codepitch*#1+0.14});
  \node[anchor=west, font=\footnotesize, text=AccentOrange, align=left] at ([xshift=0.65cm]codeR |- 0,{-\codepitch*#1+0.14}) {#2};}
\newcommand{\kw}[1]{\textcolor{VSBlue}{\textbf{#1}}}
\newcommand{\fn}[1]{\textcolor{VSPurple}{#1}}
\newcommand{\cm}[1]{\textcolor{VSComment}{\textit{#1}}}
\newcommand{\st}[1]{\textcolor{VSString}{#1}}
\newcommand{\tb}{\phantom{xxxx}}


% the generator: a dark box with cogs, like the block cipher; \gennode{name}{position}{size: 1 large, 0 small}
\newcommand{\gennode}[3]{%
  \ifnum#3=1
    \node[hwmain, minimum width=2.4cm, minimum height=1.5cm] (#1) at (#2) {};
    \node[text=white, font=\LARGE] at ([yshift=0.22cm]#1.center) {\faCogs};
    \node[text=white, font=\small\bfseries] at ([yshift=-0.42cm]#1.center) {generator};
  \else
    \node[hwmain, minimum width=1.9cm, minimum height=1.05cm] (#1) at (#2) {};
    \node[text=white, font=\large] at ([yshift=0.16cm]#1.center) {\faCogs};
    \node[text=white, font=\scriptsize\bfseries] at ([yshift=-0.3cm]#1.center) {generator};
  \fi}
\tikzset{
  bigxor/.style={circle, draw=nInk, very thick, fill=white, inner sep=0pt, minimum size=0.8cm, font=\Large},
  smallbyte/.style={minimum width=0.5cm, minimum height=0.45cm},
}







\title{Cryptography}
\subtitle{Salsa20 and ChaCha20}
\author{Ksawery Możdżyński}
\date{}

\begin{document}

% ======================================================================
% Colour mapping for the 4 by 4 state (added in Lecture 2)
%   constants : neutral (nFill), the label says "const"
%   key       : S4 amber, light fill (as every key)
%   nonce     : white with a full amber frame (public, as the IV in Lecture 1)
%   counter   : white with a dashed amber frame (public, changes by itself)
% ======================================================================
\tikzset{
  wcell/.style={draw=nLine, thick, fill=white, text=nInk, font=\ttfamily\scriptsize, minimum width=1.55cm,
    minimum height=0.55cm, inner sep=1pt},
  st-c/.style={wcell, fill=nFill},
  st-k/.style={wcell, fill=sAmber!25, draw=sAmber},
  st-n/.style={wcell, draw=sAmber, line width=1pt},
  st-t/.style={wcell, draw=sAmber, line width=1pt, dashed},
  st-x/.style={wcell},
}
% a cell that changed or matters now: a light orange fill and a thin orange frame (as in Lecture 1)
\tikzset{cellmark/.style={draw=AccentOrange, thick, fill=AccentOrange, fill opacity=0.18, draw opacity=1, rounded corners=1.5pt}}
% the 4 by 4 state: \statefig[cell width]{value/kind, ... (16 entries)}{extra TikZ}; cells are named W0 ... W15
\newcommand{\statefig}[3][1.65]{%
  \begin{tikzpicture}[>=Stealth]
    \edef\statelist{#2}
    \foreach \v/\s [count=\i from 0] in \statelist {
      \pgfmathtruncatemacro{\r}{\i/4}\pgfmathtruncatemacro{\c}{mod(\i,4)}
      \node[st-\s, minimum width=#1cm] (W\i) at ({(#1+0.08)*\c},{-0.65*\r}) {\v};
    }
    #3
  \end{tikzpicture}}
\newcommand{\statelegend}{{\scriptsize
  \tikz[baseline=-0.6ex]\node[st-c, minimum width=0.5cm, minimum height=0.3cm] {};\ constant\quad
  \tikz[baseline=-0.6ex]\node[st-k, minimum width=0.5cm, minimum height=0.3cm] {};\ key\quad
  \tikz[baseline=-0.6ex]\node[st-t, minimum width=0.5cm, minimum height=0.3cm] {};\ counter\quad
  \tikz[baseline=-0.6ex]\node[st-n, minimum width=0.5cm, minimum height=0.3cm] {};\ nonce}}
\newcommand{\markW}[1]{\draw[cellmark] (W#1.south west) rectangle (W#1.north east);}
% a black box with three inputs (key, nonce, counter) and one block out; #1 = the name in the box
\newcommand{\threeinputs}[1]{%
  \begin{tikzpicture}[>=Stealth]
    \node[actorlight=sAmber, font=\scriptsize\bfseries, minimum width=2.3cm] (k) at (0,0.9) {\faKey\ key {\normalfont(32 bytes)}};
    \node[draw=sAmber, line width=1pt, fill=white, rounded corners=3pt, font=\scriptsize\bfseries, minimum width=2.3cm, inner sep=4pt] (n) at (0,0) {nonce {\normalfont(12 bytes)}};
    \node[draw=sAmber, line width=1pt, dashed, fill=white, rounded corners=3pt, font=\scriptsize\bfseries, minimum width=2.3cm, inner sep=4pt] (t) at (0,-0.9) {counter};
    \node[hwmain, minimum width=3.0cm, minimum height=2.3cm] (box) at (3.9,0) {};
    \node[text=white, font=\LARGE] at (3.9,0.35) {\faCogs};
    \node[text=white, font=\small\bfseries, align=center] at (3.9,-0.45) {#1};
    \draw[data] (k.east) -- (box.west |- k.east); \draw[data] (n.east) -- (box.west |- n.east); \draw[data] (t.east) -- (box.west |- t.east);
    \draw[data, line width=1.6pt] (box.east) -- (6.4,0);
    \node[ksb, minimum width=2.2cm, minimum height=0.6cm] (o) at (7.6,0) {64 bytes};
    \node[font=\scriptsize, text=nInk, anchor=north] at (o.south) {one block of keystream};
  \end{tikzpicture}}
% the ChaCha20 quarter round in a code window: #1 = the line that runs now (0 = none), #2 = width
\newcommand{\qrcode}[1]{%
  \begin{tikzpicture}
    \codebox{4}{5.2}
    \cl{1}{a += b;\ \ d \textasciicircum= a;\ \ d = rotl(d, 16)}
    \cl{2}{c += d;\ \ b \textasciicircum= c;\ \ b = rotl(b, 12)}
    \cl{3}{a += b;\ \ d \textasciicircum= a;\ \ d = rotl(d, 8)}
    \cl{4}{c += d;\ \ b \textasciicircum= c;\ \ b = rotl(b, 7)}
    \ifnum#1>0 \clmark{#1}\fi
  \end{tikzpicture}}
% the Salsa20 quarter round in the same window
\newcommand{\salsaqrcode}{%
  \begin{tikzpicture}
    \codebox{4}{5.2}
    \cl{1}{b \textasciicircum= rotl(a + d, 7)}
    \cl{2}{c \textasciicircum= rotl(b + a, 9)}
    \cl{3}{d \textasciicircum= rotl(c + b, 13)}
    \cl{4}{a \textasciicircum= rotl(d + c, 18)}
  \end{tikzpicture}}


% ======================================================================
% Pictures used on more than one slide
% ======================================================================
% Mallory flips one bit on the way from Alice to Bob:
%   \tamperfig{the 11 ciphertext bytes of "Pay Bob 100"}{the same bytes after Mallory}
% byte 8 is the "1" of "Pay Bob 100": XOR 0x08 turns it into "9"
\newcommand{\tamperfig}[2]{%
  \begin{tikzpicture}[>=Stealth]
    \node[sender, minimum width=1.9cm, minimum height=0.8cm] (al) at (0,0) {\faUser\ Alice};
    \node[attacker, minimum width=3.6cm, minimum height=0.8cm] (ma) at (5.4,0) {\faUserSecret\ Mallory flips one bit};
    \node[receiver, minimum width=1.9cm, minimum height=0.8cm] (bo) at (10.8,0) {\faUser\ Bob};
    \draw[data, line width=1.6pt] (al) -- (ma) node[midway, above, font=\small\bfseries] {C};
    \draw[data, line width=1.6pt] (ma) -- (bo) node[midway, above, font=\small\bfseries] {C'};
    \node[rowlabel, font=\small] at (2.55,-1.05) {C};
    \byterow{2.95}{-1.05}{ctb}{#1}{0.68}
    \node[rowlabel, font=\small] at (2.55,-1.75) {C'};
    \byterow{2.95}{-1.75}{ctb}{#2}{0.68}
    \draw[cellmark] ({2.95+0.68*8-0.31},-1.99) rectangle ({2.95+0.68*8+0.31},-1.51);
    \node[font=\scriptsize, text=sPinkText, anchor=west] at ({2.95+0.68*10+0.45},-1.75) {one bit changed};
    \node[rowlabel, font=\small] at (2.55,-2.6) {Bob reads};
    \byterow{2.95}{-2.6}{ptb}{P,a,y,{\textvisiblespace},B,o,b,{\textvisiblespace},9,0,0}{0.68}
    \draw[cellmark] ({2.95+0.68*8-0.31},-2.84) rectangle ({2.95+0.68*8+0.31},-2.36);
    \node[font=\scriptsize, text=nInk, anchor=west] at ({2.95+0.68*10+0.45},-2.6) {Alice wrote \texttt{100}};
  \end{tikzpicture}}

% a video player, 11 by 3 units
\newcommand{\player}{%
  \draw[nInk, fill=nInk, rounded corners=4pt] (0,0) rectangle (11,3.0);
  \fill[nInk!80!black] (0.15,0.75) rectangle (10.85,2.85);
  \node[text=white!70, font=\fontsize{30}{30}\selectfont] at (5.5,1.8) {\faPlayCircle};
  \node[font=\scriptsize, text=white!80, anchor=north west] at (0.25,2.8) {encrypted film, 2 GB, 120 minutes};
  \fill[nLine] (0.4,0.38) rectangle (10.6,0.48);
  \fill[AccentOrange] (0.4,0.38) rectangle (8.05,0.48);
  \fill[white] (8.05,0.43) circle (0.12);
  \node[font=\tiny, text=white!80, anchor=north] at (0.4,0.33) {0 min};
  \node[font=\tiny, text=white!80, anchor=north] at (10.6,0.33) {120 min};}

% a processor chip: \chip{position}{1 = with an AES unit}
\newcommand{\chip}[2]{%
  \begin{scope}[shift={(#1)}]
    \draw[nInk, fill=nInk, rounded corners=2pt] (0,0) rectangle (2.4,1.9);
    \foreach \y in {0.25,0.55,...,1.7} {\draw[nLine, thick] (-0.18,\y) -- (0,\y); \draw[nLine, thick] (2.4,\y) -- (2.58,\y);}
    \node[draw=white, fill=nInk!70, rounded corners=1pt, font=\tiny, text=white, minimum width=0.9cm] at (0.6,1.2) {CPU};
    \ifnum#2=1
      \node[draw=white, fill=nInk!70, rounded corners=1pt, font=\tiny\bfseries, text=white, minimum width=0.9cm] (aes) at (1.75,1.2) {AES};
      \draw[annot] ([shift={(-0.07,-0.07)}]aes.south west) rectangle ([shift={(0.07,0.07)}]aes.north east);
    \else
      \node[draw=white, dashed, rounded corners=1pt, font=\tiny, text=white!60, minimum width=0.9cm] at (1.75,1.2) {no AES};
    \fi
    \node[draw=white, fill=nInk!70, rounded corners=1pt, font=\tiny, text=white, minimum width=2.0cm] at (1.2,0.45) {memory};
  \end{scope}}

% one row of 8 bits: \bitrow{name}{x}{y}{bits}{style of a cell}
\newcommand{\bitrow}[5]{%
  \foreach \b [count=\k from 0] in {#4} {
    \node[#5, minimum width=0.4cm, minimum height=0.42cm, font=\ttfamily\scriptsize, inner sep=0pt] (#1\k) at ({#2+0.42*\k},#3) {\b};
  }}

% a small 4 by 4 table drawn as coloured squares: \ministate{x}{y}{1 = the start (roles), 0 = mixed}
\newcommand{\ministate}[3]{%
  \foreach \r in {0,...,3} \foreach \c in {0,...,3} {
    \ifnum#3=1
      \ifnum\r=0 \node[st-c, minimum width=0.3cm, minimum height=0.3cm, inner sep=0pt] at ({#1+0.33*\c},{#2-0.33*\r}) {};\fi
      \ifnum\r=1 \node[st-k, minimum width=0.3cm, minimum height=0.3cm, inner sep=0pt] at ({#1+0.33*\c},{#2-0.33*\r}) {};\fi
      \ifnum\r=2 \node[st-k, minimum width=0.3cm, minimum height=0.3cm, inner sep=0pt] at ({#1+0.33*\c},{#2-0.33*\r}) {};\fi
      \ifnum\r=3
        \ifnum\c=0 \node[st-t, minimum width=0.3cm, minimum height=0.3cm, inner sep=0pt] at ({#1+0.33*\c},{#2-0.33*\r}) {};
        \else \node[st-n, minimum width=0.3cm, minimum height=0.3cm, inner sep=0pt] at ({#1+0.33*\c},{#2-0.33*\r}) {};\fi
      \fi
    \else
      \node[ksb, minimum width=0.3cm, minimum height=0.3cm, inner sep=0pt, rounded corners=0pt] at ({#1+0.33*\c},{#2-0.33*\r}) {};
    \fi
  }}

% ======================================================================
% Opening
% ======================================================================

\begin{frame}
  \titlepage
  \say{Dziś ChaCha20, szyfr strumieniowy, który chroni dużą część dzisiejszych połączeń HTTPS (zwłaszcza z telefonów bez sprzętowego AES), VPN-y i SSH, oraz jego starszy brat Salsa20. Zobaczymy ChaCha20 jako odpowiedź na to, co RC4 robił źle.}
\end{frame}

\begin{frame}{Agenda For Today}
  \vfill
  \begin{center}
    \begin{tikzpicture}
      % what we will do today: one icon and one sentence per point (the section names are in the bar at the top)
      \foreach \ic/\t [count=\k from 0] in {
          {\faForward}/{Build a keystream from a key, a nonce and a counter, as ChaCha20 does.},
          {\faRandom}/{Mix it so that it looks random (ChaCha20 and its older sibling Salsa20), test it where RC4 failed, and see why it is fast on any phone.},
          {\faIcon{shield-alt}}/{Add a tag, so that any change to a message is caught.}} {
        \node[circle, fill=DeepBlue, text=white, minimum size=1.3cm, font=\LARGE] at (0,{-1.9*\k}) {\ic};
        \node[anchor=west, text width=12cm, font=\large, align=left, execute at begin node={\hyphenpenalty=10000 }] at (1.0,{-1.9*\k}) {\t};
      }
    \end{tikzpicture}
  \end{center}
  \vfill
  \say{Trzy naprawy RC4, w kolejności problemów z sekcji 1: nonce i licznik jako wejścia, keystream, który wygląda losowo, i znacznik przy wiadomości. Po drodze Salsa20 (pierwsza wersja tego samego pomysłu) i krótka sekcja o szybkości. Nazwy sekcji są na pasku u góry, a link do zadań jest przy każdym zadaniu.}
\end{frame}

% ======================================================================
\section[What went wrong with RC4]{What went wrong with RC4}
% ======================================================================

\begin{frame}{Recall: a stream cipher}
  \vfill
  \begin{center}
    % CHECKED: the bytes of Lecture 1, Task 2: RC4 key "Kangaroo", "Meet m" XOR 69 9D DA 7E FE 61 = 24 F8 BF 0A DE 0C
    \begin{tikzpicture}[>=Stealth]
      \node[actorlight=sAmber, font=\small\bfseries, minimum height=0.8cm] (k) at (0,1.0) {\faKey\ key};
      \gennode{g}{2.6,1.0}{0}
      \draw[data] (k) -- (g);
      \draw[data, line width=1.4pt] (g.east) -- (4.45,1.0);
      \byterow{4.9}{1.0}{ksb}{69,9D,DA,7E,FE,61}{0.7}
      \node[font=\small] at (9.2,1.0) {\ldots};
      \node[font=\small, text=nInk, anchor=west] at (9.6,1.0) {keystream \textbf{Z}\ \ \textcolor{gray}{(looks random)}};
      \node[font=\small\bfseries, text=nInk] at (6.65,0.5) {$\oplus$};
      \node[rowlabel, font=\small] at (4.45,0) {plaintext \textbf{P}};
      \byterow{4.9}{0}{ptb}{M,e,e,t,{\textvisiblespace},m}{0.7}
      \node[font=\small] at (9.2,0) {\ldots};
      \node[font=\scriptsize\bfseries, text=nInk] at (6.65,-0.5) {$=$};
      \node[rowlabel, font=\small] at (4.45,-1.0) {ciphertext \textbf{C}};
      \byterow{4.9}{-1.0}{ctb}{24,F8,BF,0A,DE,0C}{0.7}
      \node[font=\small] at (9.2,-1.0) {\ldots};
    \end{tikzpicture}
  \end{center}
  \vfill
  \begin{center}\Large $C = P \oplus Z$\qquad\qquad $P = C \oplus Z$\\[2pt] \small ($\oplus$ = XOR, bit by bit)\end{center}
  \vfill
  \keymsg{Encrypt and decrypt are the same XOR with the same keystream.}
  \say{Krótkie przypomnienie, bo nie zakładam, że wszyscy znają szyfr strumieniowy. Generator robi z klucza długi ciąg, szyfrujemy XOR-em bajt po bajcie, odbiorca robi to samo. (Bajty na rysunku to Task 2 z wykładu 1: RC4 z kluczem Kangaroo.)}
\end{frame}

\samedot
\begin{frame}{Recall: never the same keystream twice}
  \begin{center}\small Two messages, the same keystream Z:\end{center}
  \vfill
  \begin{center}
    \begin{tikzpicture}[>=Stealth]
      \envelope{e1}{0,0.8}{$C_1$}
      \envelope{e2}{0,-0.8}{$C_2$}
      \node[actorlight=sAmber, font=\scriptsize\bfseries] at (1.45,1.25) {Z};
      \node[actorlight=sAmber, font=\scriptsize\bfseries] at (1.45,-0.35) {Z};
      \draw[data, line width=1.6pt] (1.5,0) -- (2.6,0);
      \node[font=\normalsize, text=nInk, anchor=west] at (2.7,0) {$C_1 \oplus C_2 = (P_1 \oplus Z) \oplus (P_2 \oplus Z) = \mathbf{P_1 \oplus P_2}$};
    \end{tikzpicture}
  \end{center}
  \vfill
  \begin{center}\small The keystream cancels. A \textbf{nonce} keeps it from repeating:\\
    a public value, different for every message.\end{center}
  \vfill
  \keymsg{One keystream, one message. Never twice.}
  \say{Jeśli dwie wiadomości dostaną ten sam keystream, XOR szyfrogramów usuwa klucz. Stąd nonce: jawna wartość, która zmienia keystream przy każdej wiadomości. Wrócimy do niej w sekcji o liczniku.}
\end{frame}

\begin{frame}{RC4: a stream cipher we no longer use}
  \vfill
  \begin{columns}[c,onlytextwidth]
    \column{0.56\textwidth}
      \large \textbf{RC4}, 1987. For about 20 years it protected:\\[0.3cm]
      \normalsize
      \hspace*{0.4cm}\textcolor{nInk}{\faWifi}\ \ Wi-Fi networks (WEP)\\[0.15cm]
      \hspace*{0.4cm}\textcolor{nInk}{\faLock}\ \ browser-to-website traffic (TLS)\\[0.5cm]
      Banned in TLS in 2015. Why did it fall?\\[0.2cm]
      \hspace*{0.4cm}$\rightarrow$ three problems; ChaCha20 fixes each one.
    \column{0.4\textwidth}
      \centering
      \begin{tikzpicture}
        \node[hw, minimum width=2.1cm, minimum height=2.1cm] (w) at (0,0) {};
        \node[text=nInk, font=\fontsize{30}{32}\selectfont] at (0,0.22) {\faWifi};
        \node[font=\scriptsize\bfseries, text=nInk] at (0,-0.68) {Wi-Fi (WEP)};
        \node[hw, minimum width=2.1cm, minimum height=2.1cm] (b) at (2.5,0) {};
        \node[text=nInk, font=\fontsize{30}{32}\selectfont] at (2.5,0.22) {\faWindowMaximize};
        \node[circle, fill=nInk, text=white, font=\scriptsize, inner sep=2pt] at (3.05,-0.15) {\faLock};
        \node[font=\scriptsize\bfseries, text=nInk] at (2.5,-0.68) {browser (TLS)};
        \draw[aWrite, line width=2.4pt] (-1.2,-1.2) -- (3.7,1.2);
        \node[font=\scriptsize\bfseries, text=aWrite, anchor=north] at (1.25,-1.25) {banned in TLS, 2015};
      \end{tikzpicture}
  \end{columns}
  \vfill
  \keymsg{A cipher can be everywhere and still be the wrong choice.}
  % CHECKED: RFC 7465 (February 2015) prohibits RC4 cipher suites in TLS (local/audyt/fakty.md, 2026-11-19)
  \say{RC4 był prosty i szybki, dlatego był wszędzie. Dziś jest zakazany (RFC 7465, 2015). Zaraz zobaczymy trzy problemy, ponumerowane tak samo jak naprawy dalej (pomarańczowa etykieta ,,RC4 fix N/3'' w prawym górnym rogu). Źródło zakazu: RFC 7465 (luty 2015).}
\end{frame}

\begin{frame}{RC4 problem 1: the key is its only input}
  \vfill
  \begin{center}
    \begin{tikzpicture}[>=Stealth]
      \node[font=\small\bfseries, text=nInk, anchor=west] at (-0.9,1.45) {Modern cipher};
      \node[actorlight=sAmber, font=\scriptsize\bfseries] (k1) at (0,0.8) {\faKey\ key};
      \node[draw=sAmber, line width=1pt, fill=white, rounded corners=3pt, font=\scriptsize\bfseries, inner sep=4pt] (n1) at (0,0.05) {nonce};
      \node[draw=sAmber, line width=1pt, dashed, fill=white, rounded corners=3pt, font=\scriptsize\bfseries, inner sep=4pt] (t1) at (0,-0.7) {counter};
      \node[hwmain, minimum width=1.9cm, minimum height=1.9cm, font=\small\bfseries] (c1) at (2.6,0.05) {cipher};
      \draw[data] (k1) -- (c1.west |- k1); \draw[data] (n1) -- (c1.west |- n1); \draw[data] (t1) -- (c1.west |- t1);
      \draw[data, line width=1.4pt] (c1.east) -- (4.3,0.05);
      \byterow{4.7}{0.05}{ksb}{,,,}{0.55}
      \node[font=\scriptsize, text=nInk, anchor=north] at (5.5,-0.25) {keystream};
      \node[font=\small\bfseries, text=nInk, anchor=west] at (7.3,1.45) {RC4};
      \node[actorlight=sAmber, font=\scriptsize\bfseries] (k2) at (8.2,0.8) {\faKey\ key};
      \node[hwmain, minimum width=1.9cm, minimum height=1.9cm, font=\small\bfseries] (c2) at (10.8,0.05) {\faCogs\ RC4};
      \draw[data] (k2) -- (c2.west |- k2);
      \draw[annot, dashed] (7.5,-0.95) rectangle (8.9,0.3);
      \node[font=\scriptsize, text=AccentOrange, align=center] at (8.2,-0.33) {no other\\ input};
      \draw[data, line width=1.4pt] (c2.south) -- (10.8,-1.25);
      \byterow{10.0}{-1.55}{ksb}{,,,}{0.55}
      \node[font=\scriptsize, text=nInk, anchor=north] at (10.8,-1.85) {keystream, from byte 0 on};
    \end{tikzpicture}
  \end{center}
  \vfill
  \begin{center}\small No nonce: everyone glues it in their own way. WEP glued it to the key, and it broke.\\[3pt]
    No counter: RC4 only goes forward. To get byte $n$, it must make all the bytes before it.\end{center}
  \vfill
  \keymsg{One input only: no safe place for a nonce, no way to jump to the middle.}
  \say{RC4 bierze tylko klucz. Nonce trzeba dokleić samemu, a to prosta droga do błędu: WEP dokleił go do klucza i dał się złamać. Do tego RC4 idzie tylko do przodu: żeby dostać bajt numer n, trzeba wyprodukować wszystkie wcześniejsze. Dobry szyfr ma nonce i licznik (numer bloku) jako osobne wejścia. To jest problem 1 i naprawa 1 w sekcji 2.}
\end{frame}

\begin{frame}{RC4 problem 2: the keystream does not look random}
  \begin{center}\small Two symptoms of \textbf{one} weakness:\end{center}
  \vfill
  \begin{columns}[T,onlytextwidth]
    \column{0.48\textwidth}
      \centering\textbf{(a) the output remembers the key}\\[0.2cm]
      \begin{tikzpicture}[>=Stealth]
        \node[actorlight=sAmber, font=\scriptsize\bfseries, minimum height=0.7cm] (in) at (0,0) {IN: \faKey\ key};
        \node[hwmain, minimum width=1.5cm, minimum height=0.8cm, font=\small\bfseries] (r) at (2.4,0) {\faCogs\ RC4};
        \node[ksb, minimum width=1.3cm, minimum height=0.7cm, font=\scriptsize\bfseries] (out) at (4.6,0) {OUT};
        \draw[data] (in) -- (r); \draw[data] (r) -- (out);
        \draw[->, very thick, AccentOrange, dashed] (out.north) .. controls (4.6,1.3) and (0,1.3) .. (in.north);
        \node[font=\scriptsize\bfseries, text=AccentOrange, fill=white, inner sep=1pt] at (2.3,1.0) {remembers};
      \end{tikzpicture}\\[0.2cm]
      \small OUT still ``remembers'' IN\\[0.1cm]
      $\rightarrow$ WEP key recovered from traffic\\[0.5cm]
      {\scriptsize\textcolor{gray}{\faSearch\ For the curious:}\\ \simlink{rc4/\#4.1}\\ \simlink{rc4/\#4.2}}
    \column{0.48\textwidth}
      \centering\textbf{(b) the stream is biased}\\[0.2cm]
      % CHECKED: Lecture 1, checks/bias_10M.c: the 2nd RC4 byte is 0 in 0.784% of keys, any value: 0.391%
      \begin{tikzpicture}
        \foreach \lab/\v [count=\k from 0] in {0/0.784, 1/0.391, 2/0.391, 3/0.391, 4/0.391, {\ldots}/0, 255/0.391} {
          \ifnum\k=0 \fill[AccentOrange] ({0.55*\k-0.18},0) rectangle ({0.55*\k+0.18},{2.0*\v});
          \else \fill[nLine] ({0.55*\k-0.18},0) rectangle ({0.55*\k+0.18},{2.0*\v}); \fi
          \node[font=\tiny, text=nInk, anchor=north] at ({0.55*\k},0) {\lab};
        }
        \draw[nLine] (-0.35,0) -- (3.65,0);
        \node[font=\tiny\bfseries, text=AccentOrange, anchor=south] at (0,1.57) {0.78\%};
        \node[font=\tiny, text=nInk, anchor=south] at (1.65,0.78) {0.39\%};
        \node[font=\tiny, text=gray, anchor=north] at (1.65,-0.22) {value of the 2nd keystream byte};
      \end{tikzpicture}\\[0.1cm]
      \small second keystream byte is 0\\ twice as often as it should\\[0.1cm]
      $\rightarrow$ a repeated message leaks\\ (e.g.\ a browser cookie)
  \end{columns}
  \vfill
  \keymsg{A good keystream looks random. RC4's does not, in two ways.}
  \say{Problem 2 ma dwa objawy, ale to jedna słabość: keystream nie wygląda losowo. Raz widać to od strony klucza (wyjście go pamięta), raz od strony wyjścia (drugi bajt jest skrzywiony). Liczby pokażę w sekcji 3, na żywo (naprawa 2). Źródła: ciasteczko z powtarzanej wiadomości to AlFardan i in. (2013) oraz RC4 NOMORE (Vanhoef, Piessens, 2015). Link na dole: symulator RC4 z wykładu 1, gdyby ktoś chciał to sobie przypomnieć.}
\end{frame}

\begin{frame}{RC4 problem 3: no integrity}
  \vfill
  \begin{center}
    % CHECKED: python3, RC4 key "Kangaroo", "Pay Bob 100" XOR the keystream = the first row;
    %          byte 8 XOR 0x08 (3D -> 35) decrypts to "Pay Bob 900"
    \tamperfig{39,FC,A3,5E,BC,0E,11,42,3D,70,2A}{39,FC,A3,5E,BC,0E,11,42,35,70,2A}
  \end{center}
  \vfill
  \begin{center}\small Bob decrypts C' and reads a changed message. He has no way to know.\\[3pt]
    A stream cipher hides the message. It does \textbf{not} protect it from change.\end{center}
  \vfill
  \defbox{\textbf{Integrity}: the certainty that nobody changed the message.}
  \keymsg{Three fixes to come: 1 a nonce and a counter as inputs, 2 a keystream that looks random, 3 a tag.}
  \say{To nie jest błąd tylko RC4: żaden goły szyfr strumieniowy nie chroni integralności. Mallory odwraca bit w szyfrogramie, a odbiorca niczego nie zauważa. (Na rysunku RC4 z kluczem Kangaroo: bajt 3D zmienia się w 35, a 100 w 900.) Podsumowanie: trzy naprawy, które zobaczymy po kolei, w tej samej kolejności co problemy.}
\end{frame}

% ======================================================================
\section[A keystream from a counter]{A keystream from a counter}
% ======================================================================

\begin{frame}{A film on a server}
  \begin{center}\small An encrypted film, 2 GB. The viewer jumps to minute 90.\end{center}
  \vspace{-0.1cm}
  \begin{center}
    \begin{tikzpicture}[>=Stealth]
      \begin{scope}[scale=0.62, transform shape]
        \player
      \end{scope}
      % the keystream tape under the player: 0 to 120 min over 6.82 units, minute 90 at 5.0
      \begin{scope}[shift={(0,-0.9)}]
        \foreach \k/\l in {0/b0, 1/b1, 2/b2} \node[ksb, fill=nIdle, draw=nLine, minimum width=0.5cm, minimum height=0.45cm, font=\ttfamily\tiny] at ({0.27+0.52*\k},0) {\l};
        \fill[nIdle] (1.55,-0.225) rectangle (4.75,0.225);
        \draw[nLine, thick] (1.55,-0.225) rectangle (4.75,0.225);
        \node[font=\scriptsize, text=gray] at (3.15,0) {\ldots};
        \node[ksb, minimum width=0.5cm, minimum height=0.45cm, font=\ttfamily\tiny] (need) at (5.0,0) {b?};
        \draw[nLine, thick] (5.25,-0.225) rectangle (6.82,0.225);
        \draw[->, very thick, AccentOrange] (7.6,0) -- (need.east);
        \node[font=\scriptsize\bfseries, text=AccentOrange, anchor=west] at (7.65,0) {we need this part};
        \draw[decorate, decoration={brace, mirror, amplitude=5pt}, nInk, thick] (0.02,-0.32) -- (4.75,-0.32)
          node[midway, below=6pt, font=\scriptsize, text=nInk] {RC4 must make and throw away all of this first};
        \node[font=\scriptsize, text=nInk, anchor=east] at (-0.1,0) {keystream};
      \end{scope}
    \end{tikzpicture}
  \end{center}
  \vspace{-0.1cm}
  \begin{center}\small RC4 only goes forward, byte by byte from the start.\end{center}
  \vfill
  \fixtag{1}%
  \keymsg{We want any part of the keystream, without the parts before it.}
  \say{Naprawa 1, druga połowa problemu 1 (,,RC4 idzie tylko do przodu''). Przewijanie filmu: chcemy keystream w środku, nie od początku. RC4 jako maszyna idąca tylko do przodu musiałby policzyć i wyrzucić wszystko wcześniej, tu 1,5 GB. Potrzebujemy czegoś lepszego.}
\end{frame}

\begin{frame}{The idea: number the blocks}
  \begin{center}\small Cut the keystream into \textbf{64-byte blocks}. Number them 0, 1, 2, \ldots\end{center}
  \vfill
  \begin{center}
    \large block $n$ = f(\,\tikz[baseline=-0.6ex]\node[actorlight=sAmber, font=\small\bfseries, inner sep=2pt] {key};\,,\
    \tikz[baseline=-0.6ex]\node[draw=sAmber, line width=1pt, rounded corners=3pt, font=\small\bfseries, inner sep=2pt] {nonce};\,,\
    \tikz[baseline=-0.6ex]\node[draw=sAmber, line width=1pt, dashed, rounded corners=3pt, font=\small\bfseries, inner sep=2pt] {$n$};\,)
  \end{center}
  \vfill
  \begin{center}
    \begin{tikzpicture}[>=Stealth]
      \foreach \n/\x in {0/0, 1/1.9, 2/3.8, n/7.2} {
        \node[hwmain, minimum width=0.9cm, minimum height=0.5cm, font=\scriptsize\bfseries, inner sep=1pt] (f\n) at (\x,1.05) {f};
        \node[font=\tiny, text=nInk, anchor=south] at (\x,1.32) {key, nonce, \n};
        \node[ksb, minimum width=1.75cm, minimum height=0.6cm] (b\n) at (\x,0) {block \n};
        \draw[data] (f\n.south) -- (b\n.north);
      }
      \node[font=\large, text=nInk] at (5.5,0) {\ldots};
      \node[font=\scriptsize, text=nInk, anchor=west] at (8.3,0.5) {each block made on its own};
    \end{tikzpicture}
  \end{center}
  \vfill
  \bigdef{\textbf{Counter}: a block number, which 64-byte block of the keystream we want.}
  \fixtag{1}%
  \keymsg{Block $n$ depends only on the key, the nonce and $n$.}
  \say{Zamiast maszyny idącej do przodu: każdy 64-bajtowy blok liczymy osobno z klucza, nonce i numeru bloku. Numer bloku to licznik. Dzięki temu dojdziemy od razu do bloku 700.}
\end{frame}

\samedot
\begin{frame}{Why separate blocks are safe}
  \vfill
  \begin{center}
    \begin{tikzpicture}[>=Stealth]
      \foreach \lab/\y/\shades in {{key, N1, 0}/1.2/{15,70,35,90,10,55,80,25,60,40}, {key, N1, 1}/0.5/{85,20,60,30,75,10,45,95,15,65}, {key, N2, 0}/-0.2/{40,95,10,65,25,85,50,5,70,30}} {
        \node[font=\scriptsize\ttfamily, text=nInk, anchor=east] at (1.6,\y) {\lab};
        \node[hwmain, minimum width=0.9cm, minimum height=0.45cm, font=\scriptsize\bfseries, inner sep=1pt] (f) at (2.3,\y) {f};
        \draw[data] (f.east) -- (3.2,\y);
        \foreach \sh [count=\c from 0] in \shades \node[draw=sAmber, fill=sAmber!\sh, minimum width=0.3cm, minimum height=0.4cm, inner sep=0pt] at ({3.45+0.32*\c},\y) {};
      }
      \node[font=\scriptsize, text=nInk, anchor=west] at (6.9,0.5) {look unrelated and random};
      \node[font=\scriptsize\ttfamily, text=nInk, anchor=east] at (1.6,-1.2) {key, N1, 0};
      \node[hwmain, minimum width=0.9cm, minimum height=0.45cm, font=\scriptsize\bfseries, inner sep=1pt] (g) at (2.3,-1.2) {f};
      \draw[data] (g.east) -- (3.2,-1.2);
      \foreach \sh [count=\c from 0] in {15,70,35,90,10,55,80,25,60,40} \node[draw=sAmber, fill=sAmber!\sh, minimum width=0.3cm, minimum height=0.4cm, inner sep=0pt] at ({3.45+0.32*\c},-1.2) {};
      \node[font=\scriptsize\bfseries, text=AccentOrange, anchor=west] at (6.9,-1.2) {the same input again: the same block as the first row};
      \draw[nLight, thick, dashed] (-0.3,-0.65) -- (12.6,-0.65);
    \end{tikzpicture}
  \end{center}
  \vfill
  \defbox{\textbf{Pseudorandom function}: without the key, its outputs look like independent random blocks, for every different input.}
  \begin{center}\small A new (nonce, $n$) gives a block that looks fresh: all blocks together look like one long random keystream.\\[3pt]
    The same (nonce, $n$) gives the same block again. So this pair must never repeat under one key.\end{center}
  \vfill
  \fixtag{1}%
  \keymsg{Each block looks fresh and random, as long as its input is new.}
  \say{Dlaczego wolno liczyć każdy blok osobno: dobre f zachowuje się jak funkcja losowa. Kto nie zna klucza, nie odróżni wyników dla różnych wejść od niezależnych losowych bloków, nawet jeśli wejścia różnią się o jeden bit (licznik 0 i 1). Dlatego bloki 0, 1, 2, ... razem wyglądają jak jeden długi losowy keystream. Jedyny wyjątek: to samo wejście daje ten sam blok, stąd reguła, że para (nonce, licznik) nie może się powtórzyć. Ten sam pomysł zobaczymy w sekcji 4 przy AES-CTR.}
\end{frame}

% CHECKED: Salsa20 2005, ChaCha20 2008, Bernstein: https://cr.yp.to/chacha.html (local/audyt/fakty.md)
\begin{frame}{ChaCha20 from the outside}
  \vfill
  \begin{center}\threeinputs{ChaCha20}\end{center}
  \vfill
  \begin{columns}[c,onlytextwidth]
    \column{0.45\textwidth}
      \centering
      {\scriptsize RC4, for contrast:}\\[3pt]
      \begin{tikzpicture}[>=Stealth]
        \node[actorlight=sAmber, font=\scriptsize\bfseries] (k) at (0,0) {\faKey\ key};
        \node[hwmain, minimum width=1.3cm, minimum height=0.55cm, font=\scriptsize\bfseries] (r) at (1.9,0) {RC4};
        \draw[data] (k) -- (r);
        \draw[data] (r) -- (3.3,0);
        \node[font=\scriptsize, text=nInk, anchor=west] at (3.35,0) {keystream};
        \node[font=\scriptsize, text=gray, anchor=north] at (1.9,-0.35) {(one input)};
      \end{tikzpicture}
    \column{0.52\textwidth}
      \small Daniel J. Bernstein: \textbf{Salsa20} (2005),\\ then \textbf{ChaCha20} (2008), which mixes faster\\ (we compare them in Section 3).\\[3pt]
      Today ChaCha20 is the usual choice.
  \end{columns}
  \vfill
  \curious{chacha20/\#1}
  \fixtag{1}%
  \keymsg{Three inputs in, 64 bytes out. We use it before we open it.}
  \say{ChaCha20 jako czarna skrzynka: klucz, nonce, licznik na wejściu, 64 bajty na wyjściu. Od razu widać różnicę z RC4: nonce ma własne wejście. Najpierw go użyjemy, dopiero potem zajrzymy do środka. W symulatorze z linku można zmieniać licznik i patrzeć, jak wiadomość przesuwa się po keystreamie. Gdyby ktoś pytał, skąd Alice i Bob mają wspólny 32-bajtowy klucz: to wymiana kluczy z wykładów 4 i 5 (Diffie-Hellman, krzywe eliptyczne).}
\end{frame}

\begin{frame}{Task 1: use ChaCha20}
  \vfill
  \taskgoal{decrypt a message from RFC 8439 with one call to ChaCha20.}
  \taskready{\texttt{chacha.py}, the function \texttt{chacha20(key, nonce, counter, data)}.}
  \vfill
  \opentask{Task1}{about 4 minutes}
  \vfill
  \fixtag{1}%
  \say{Około 4 minut. chacha.py jest gotowy, studenci używają go jako czarnej skrzynki. Instrukcja i TODO są w README na GitHubie. Deszyfrowanie to to samo wywołanie co szyfrowanie: XOR z tym samym keystreamem. Dane z RFC 8439, punkt 2.4.2.}
\end{frame}

\samedot
\begin{frame}{Task 1: what we saw}
  \vfill
  % CHECKED: Task1/solution/use_chacha.py prints the plaintext of RFC 8439 2.4.2;
  %          chacha.py gives the ciphertext 6e2e359a2568f980... of the RFC
  \begin{center}
    \begin{minipage}{0.72\textwidth}
      \stepbox{Output}{\texttt{Ladies and Gentlemen of the class of '99: If I could offer you only one tip for the future, sunscreen would be it.}}
    \end{minipage}\\[0.4cm]
    Decrypting is the same call as encrypting, with counter 1.
  \end{center}
  \vfill
  \fixtag{1}%
  \keymsg{One call decrypts. Next: decrypt only the middle, with the counter.}
  \say{Sprawdzone: rozwiązanie daje tekst z RFC. Jedno wywołanie chacha20 z licznikiem 1 i tym samym kluczem oraz nonce. Dolne zdanie prowadzi do zadania 2: skok w środek przez licznik.}
\end{frame}

\begin{frame}{Task 2: jump to the middle}
  \vfill
  \taskgoal{decrypt \textbf{only} block 700 of an encrypted film, using the counter.}
  \taskready{\texttt{chacha.py} and \texttt{film.bin} (1000 blocks of 64 bytes).}
  \vfill
  \opentask{Task2}{about 5 minutes}
  \vfill
  \fixtag{1}%
  \say{Około 5 minut. Pokazać, że licznik to numer bloku: ustaw counter = 700 i odszyfruj sam ten kawałek. Ostatni TODO to pomiar: jeden blok kontra 701. Klucz i nonce są w pliku, instrukcja w README.}
\end{frame}

\samedot
\begin{frame}{Task 2: what we saw}
  \vfill
  % CHECKED: Task2/solution/jump.py, author's laptop (i5-8265U, Python 3.12), 2026-11-19, five runs:
  %          one block 0.00014 to 0.00020 s, all 701 blocks 0.102 to 0.114 s; shown: the first run
  %          (0.00014 s and 0.10195 s, the same run as in Task2/README.md); 0.10195 / 0.00014 = 728
  \begin{center}
    \begin{minipage}{0.6\textwidth}
      \stepbox{Output}{\texttt{Frame 0700: the hero finally opens the red door.}\\[3pt]
        \texttt{one block: 0.00014 s}\\
        \texttt{all 701:\ \ \ 0.10 s}}
    \end{minipage}\\[0.5cm]
    \begin{tikzpicture}
      \fill[nLine] (0,0.5) rectangle (6.4,0.8);
      \node[font=\small, text=nInk, anchor=west] at (6.55,0.65) {all 701 blocks: 0.10 s};
      \fill[AccentOrange] (0,-0.15) rectangle (0.06,0.15);
      \node[font=\small, text=nInk, anchor=west] at (0.2,0) {one block: about 700 times faster};
    \end{tikzpicture}
  \end{center}
  \vfill
  \fixtag{1}%
  \keymsg{The counter gives random access to the keystream.}
  \say{Jeden blok zamiast 701. To jest dostęp swobodny do keystreamu, czego RC4 nie miał. Czasy z mojego laptopa (i5-8265U, Python 3.12): pięć uruchomień dało 0,00014 do 0,00020 s dla jednego bloku i 0,10 do 0,11 s dla 701 bloków; na slajdzie i w README jest ten sam, pierwszy pomiar. U studentów liczby będą inne, ale różnica podobna.}
\end{frame}

\begin{frame}{Nonce and counter: two different jobs}
  \vfill
  \begin{center}
    \large
    \begin{tabular}{l l}
      \tikz[baseline=-0.6ex]\node[draw=sAmber, line width=1pt, rounded corners=3pt, font=\small\bfseries, inner sep=3pt] {nonce}; & changes between \textbf{messages} \\[6pt]
      \tikz[baseline=-0.6ex]\node[draw=sAmber, line width=1pt, dashed, rounded corners=3pt, font=\small\bfseries, inner sep=3pt] {counter}; & changes between \textbf{blocks} of one message \\
    \end{tabular}
  \end{center}
  \vfill
  \begin{center}
    \begin{tikzpicture}
      \foreach \m/\nn/\y in {1/N1/0.45, 2/N2/-0.45} {
        \node[font=\small\bfseries, text=nInk, anchor=east] at (-0.2,\y) {message \m:};
        \node[draw=sAmber, line width=1pt, fill=white, rounded corners=2pt, font=\scriptsize\bfseries, minimum width=1.3cm, minimum height=0.5cm] at (0.6,\y) {nonce \nn};
        \foreach \c in {0,1,2} \node[draw=sAmber, line width=1pt, dashed, fill=sAmber!25, rounded corners=1.5pt, font=\ttfamily\scriptsize, minimum width=1.2cm, minimum height=0.5cm] at ({2.1+1.3*\c},\y) {ctr \c};
        \node[font=\small, text=nInk] at (5.8,\y) {\ldots};
      }
    \end{tikzpicture}
  \end{center}
  \vfill
  \begin{center}\small The pair (nonce, counter) must \textbf{never repeat} under one key.\end{center}
  \vfill
  \fixtag{1}%
  \keymsg{The nonce is an input, not something you glue to the key.}
  \say{Nonce rozróżnia wiadomości, licznik bloki w jednej wiadomości. To domyka naprawę: nonce jest wejściem szyfru. W RC4 trzeba go było doklejać, w ChaCha20 po prostu się go podaje.}
\end{frame}

\samedot
\begin{frame}{How long can one message be?}
  \vfill
  \begin{center}
    \begin{tikzpicture}[>=Stealth]
      \foreach \k in {0,...,31} \node[draw=sAmber, dashed, thick, fill=white, minimum width=0.26cm, minimum height=0.4cm, inner sep=0pt] at ({0.28*\k},0) {};
      \node[font=\scriptsize\bfseries, text=nInk, anchor=south] at (4.34,0.3) {counter is 32-bit: $2^{32}$ blocks};
      \draw[data] (9.2,0) -- (10.2,0);
      \node[font=\small, text=nInk, anchor=west] at (10.3,0) {\textbf{256 GiB}};
    \end{tikzpicture}
  \end{center}
  \vfill
  \begin{center}\large $2^{32}$ blocks $\times$ 64 bytes = \textbf{256 GiB}\\[3pt] \normalsize per one (key, nonce) pair.\end{center}
  \vfill
  \fixtag{1}%
  \keymsg{Plenty for a message; change the nonce for the next one.}
  % CHECKED: 2**32 * 64 / 2**30 = 256
  \say{256 GiB na parę klucz i nonce. Więcej niż każda rozsądna wiadomość; dla kolejnej bierzemy nowy nonce.}
\end{frame}

% ======================================================================
\section[The keystream looks random]{The keystream looks random}
% ======================================================================

\begin{frame}{Inside: a 4 by 4 table}
  \begin{center}
    \statefig[1.45]{const/c,const/c,const/c,const/c,key/k,key/k,key/k,key/k,key/k,key/k,key/k,key/k,ctr/t,nonce/n,nonce/n,nonce/n}{%
      \node[font=\small, text=nInk, anchor=west] at ([xshift=0.3cm]W3.east) {constants (fixed)};
      \node[font=\small, text=sAmberText, anchor=west] at ([xshift=0.3cm]W7.east) {your key};
      \node[font=\small, text=sAmberText, anchor=west] at ([xshift=0.3cm]W11.east) {your key};
      \node[font=\small, text=nInk, anchor=west] at ([xshift=0.3cm]W15.east) {counter + nonce};}\\[2pt]
    \statelegend
  \end{center}
  \vfill
  \defbox{\centering \textbf{State}: everything the cipher works on, 16 words, 4 by 4.\\[1pt]
    {\footnotesize \textbf{word} = 32 bits = 4 bytes, lowest byte first (little-endian); key = 8 words, counter = 1, nonce = 3.}}
  \curious{chacha20/\#2}
  \fixtag{2}%
  \keymsg{Everything that should change the keystream is in the table.}
  \say{Wnętrze ChaCha20 to jedna tabela 4 na 4: stałe, klucz, licznik, nonce. To wszystko, na czym szyfr pracuje. Nie uczymy się liczb, tylko co gdzie siedzi. (Słowo ma 32 bity, czyli 4 bajty: klucz to 8 słów, nonce 3, licznik 1. Bajty wchodzą do słowa od najmłodszego: stała ,,expa'' to bajty 65 78 70 61, czyli słowo 61707865.) W symulatorze z linku widać tę tabelę dla własnego klucza i nonce.}
\end{frame}

\begin{frame}{Three simple steps, no tables}
  \vfill
  \begin{center}
    % CHECKED: python3, on 8 bits: 0x35 + 0xE0 = 0x115 -> 0x15 (00010101); rotl(0x35, 2) = 0xD4 (11010100);
    %          0x35 ^ 0x5C = 0x69 (01101001)
    \begin{tikzpicture}[>=Stealth]
      % add
      \node[font=\normalsize\bfseries, text=DeepBlue, anchor=west] at (-0.2,1.55) {add};
      \node[font=\scriptsize\ttfamily, text=nInk, anchor=west] at (0.6,1.55) {a = a + b};
      \node[font=\tiny, text=gray, anchor=west] at (-0.2,1.15) {(wrap at 32 bits; here on 8 bits)};
      \bitrow{A}{0.2}{0.6}{0,0,1,1,0,1,0,1}{wcell}
      \bitrow{B}{0.2}{0.1}{1,1,1,0,0,0,0,0}{wcell}
      \node[font=\small, text=nInk] at (-0.2,0.1) {+};
      \draw[nInk, thick] (-0.05,-0.2) -- (3.35,-0.2);
      \bitrow{S}{0.2}{-0.5}{0,0,0,1,0,1,0,1}{st-k}
      \node[font=\ttfamily\scriptsize, text=aWrite] (cy) at (-0.3,-0.5) {1};
      \draw[aWrite, thick] (-0.45,-0.65) -- (-0.15,-0.35);
      \node[font=\tiny, text=aWrite, anchor=north] at (-0.3,-0.75) {lost};
      % rotate
      \begin{scope}[shift={(4.7,0)}]
        \node[font=\normalsize\bfseries, text=DeepBlue, anchor=west] at (-0.2,1.55) {rotate};
        \node[font=\scriptsize\ttfamily, text=nInk, anchor=west] at (1.05,1.55) {a = a <\kern-1pt<\kern-1pt< 2};
        \node[font=\tiny, text=gray, anchor=west] at (-0.2,1.15) {bits move left, wrap (here 8 bits, by 2)};
        \bitrow{R}{0.2}{0.6}{0,0,1,1,0,1,0,1}{wcell}
        \bitrow{T}{0.2}{-0.5}{1,1,0,1,0,1,0,0}{st-k}
        \foreach \k in {2,...,7} {\pgfmathtruncatemacro{\j}{\k-2} \draw[data, thin, nLine] (R\k.south) -- (T\j.north);}
        \draw[->, thick, AccentOrange] (R0.south) .. controls (0.2,0.0) and (2.9,0.1) .. (T6.north);
        \draw[->, thick, AccentOrange] (R1.south) .. controls (0.6,-0.05) and (3.3,0.05) .. (T7.north);
      \end{scope}
      % xor
      \begin{scope}[shift={(9.4,0)}]
        \node[font=\normalsize\bfseries, text=DeepBlue, anchor=west] at (-0.2,1.55) {xor};
        \node[font=\scriptsize\ttfamily, text=nInk, anchor=west] at (0.55,1.55) {a = a \textasciicircum{} b};
        \node[font=\tiny, text=gray, anchor=west] at (-0.2,1.15) {1 where the bits differ};
        \bitrow{X}{0.2}{0.6}{0,0,1,1,0,1,0,1}{wcell}
        \bitrow{Y}{0.2}{0.1}{0,1,0,1,1,1,0,0}{wcell}
        \node[font=\small, text=nInk] at (-0.2,0.1) {$\oplus$};
        \draw[nInk, thick] (-0.05,-0.2) -- (3.35,-0.2);
        \bitrow{Z}{0.2}{-0.5}{0,1,1,0,1,0,0,1}{st-k}
      \end{scope}
    \end{tikzpicture}
  \end{center}
  \vfill
  \begin{center}\small = \textbf{ARX} (add, rotate, xor). No lookup tables in memory. ChaCha20 rotates by 16, 12, 8 and 7.\\[3pt]
    RC4 reads a table (S); ChaCha20 does not.\end{center}
  \vfill
  \curious{chacha20/\#4}
  \fixtag{2}%
  \keymsg{Three cheap steps on words. Nothing depends on a memory table.}
  \say{Cały szyfr to trzy tanie operacje na słowach: dodawanie, rotacja, XOR. Nazywa się to ARX. Na rysunku 8 bitów, żeby było widać; w ChaCha20 słowa mają 32 bity, a rotacje są o 16, 12, 8 i 7. Ważne: żadnej tablicy w pamięci, inaczej niż RC4. Wrócę do tego w sekcji o szybkości. Link: te trzy operacje bit po bicie, na 32-bitowych słowach.}
\end{frame}

\begin{frame}{Mix it twenty times}
  \vfill
  \begin{center}
    % CHECKED: RFC 8439 2.3.2: word 0 of the start 61707865, of the output block e4e7f110
    \begin{tikzpicture}[>=Stealth]
      \ministate{0}{0.5}{1}
      \node[font=\scriptsize\bfseries, text=nInk, anchor=south, align=center] at (0.5,0.75) {state\\[-2pt]{\tiny\ttfamily\mdseries 61707865\ldots}};
      \foreach \k/\x/\l in {1/2.3/1, 2/3.6/2, 20/6.2/20} {
        \node[hwmain, minimum width=1.0cm, minimum height=0.75cm, font=\scriptsize\bfseries] (m\k) at (\x,0) {mix};
        \node[font=\tiny, text=nInk, anchor=north] at (\x,-0.42) {\l};
      }
      \node[font=\normalsize, text=nInk] (dots) at (4.9,0) {\ldots};
      \draw[data] (1.25,0) -- (m1.west); \draw[data] (m1) -- (m2); \draw[data] (m2) -- (dots); \draw[data] (dots) -- (m20);
      \node[circle, draw=AccentOrange, very thick, fill=white, font=\large\bfseries, text=AccentOrange, inner sep=1pt, minimum size=0.6cm] (plus) at (7.6,0) {+};
      \draw[data] (m20) -- (plus);
      \draw[->, very thick, AccentOrange, dashed] (0.5,-0.85) -- (0.5,-1.25) -- (7.6,-1.25) -- (plus.south);
      \node[font=\scriptsize\bfseries, text=AccentOrange, fill=white, inner sep=1pt] at (4.0,-1.25) {+ initial state};
      \draw[data] (plus) -- (8.6,0);
      \ministate{8.75}{0.5}{0}
      \node[font=\scriptsize\bfseries, text=nInk, anchor=south, align=center] at (9.25,0.75) {out: 64 bytes\\[-2pt]{\tiny\ttfamily\mdseries e4e7f110\ldots}};
    \end{tikzpicture}
  \end{center}
  \vfill
  \begin{center}
    \begin{tikzpicture}
      \node[font=\scriptsize\bfseries, text=nInk, anchor=east] at (-0.2,0) {round:};
      \foreach \k in {1,...,20} {
        \pgfmathtruncatemacro{\sh}{5*\k}
        \node[draw=nLine, fill=sAmber!\sh, minimum width=0.5cm, minimum height=0.45cm, font=\tiny, text=nInk, inner sep=0pt] at ({0.55*(\k-1)},0) {\k};
      }
    \end{tikzpicture}
  \end{center}
  \vfill
  \begin{center}\small Each round stirs the whole table with add, rotate and xor, until every output bit depends on every input bit.\\[3pt]
    At the end: add the starting table back. (Without it, you could undo the rounds and read the key.)\end{center}
  \vfill
  \curious{chacha20/\#3}
  \fixtag{2}%
  \keymsg{Twenty rounds, until everything depends on everything.}
  \say{Te trzy kroki powtarzamy na całej tabeli 20 razy. Nie rysujemy ćwiartki rundy ani kolumn i przekątnych: intuicja to ,,miesza w kółko, aż wszystko zależy od wszystkiego''. Na końcu dodajemy tabelę początkową, inaczej dałoby się rundy cofnąć i odczytać klucz. Po co stałe ,,expand 32-byte k'': bez nich tabela z samych zer (klucz zero, nonce zero, licznik zero) dawałaby same zera, bo ARX zera nie rusza; stałe odróżniają też wariant z kluczem 16-bajtowym (,,expand 16-byte k''). (Liczby pod tabelami to pierwsze słowo przykładu z RFC 8439, punkt 2.3.2.) Kto chce zobaczyć ćwiartki rundy, kolumny i przekątne, ma je w symulatorze z linku.}
\end{frame}

% CHECKED: Task3/salsa.py matches the Salsa20 specification (cr.yp.to/snuffle/spec.pdf): quarterround(1,0,0,0) =
%          (08008145, 00000080, 00010200, 20500000), the Salsa20 hash example 211,159,... -> 109,42,178,168,...;
%          local/audyt/tmp-L02/work/salsa_check.py: scrypt built on this Salsa20/8 = hashlib.scrypt
% CHECKED 2026-11-25: eSTREAM final portfolio announced 15 April 2008, profile 1 (software): HC-128, Rabbit, Salsa20/12, SOSEMANUK
%   (ECRYPT II deliverable D.SYM.10, cordis.europa.eu/docs/projects/cnect/6/216676/080/deliverables/001-DSYM10v1.pdf, via search summary)
\begin{frame}{Salsa20: the older sibling}
  \vfill
  \begin{center}
    \begin{tabular}{c@{\qquad\qquad}c}
      \textbf{Salsa20} (2005) & \textbf{ChaCha20} (2008) \\[3pt]
      \statefig[1.0]{const/c,key/k,key/k,key/k,key/k,const/c,nonce/n,nonce/n,ctr/t,ctr/t,const/c,key/k,key/k,key/k,key/k,const/c}{} &
      \statefig[1.0]{const/c,const/c,const/c,const/c,key/k,key/k,key/k,key/k,key/k,key/k,key/k,key/k,ctr/t,nonce/n,nonce/n,nonce/n}{} \\[3pt]
      {\scriptsize constants on the diagonal; nonce 8 bytes, counter 8 bytes} & {\scriptsize constants in the first row; nonce 12 bytes, counter 4 bytes}
    \end{tabular}\\[6pt]
    \statelegend
  \end{center}
  \vfill
  \begin{center}\small The same idea: a 4 by 4 table, add, rotate and xor, 20 rounds, the start table added back.\\[3pt]
    In 2008 eSTREAM, a European contest for stream ciphers, chose Salsa20/12 (12 rounds) for software.\end{center}
  \vfill
  \keymsg{ChaCha20 is Salsa20 reworked: the same pieces, a new layout and a new quarter round.}
  \say{Salsa20 to pierwsza wersja tego samego pomysłu, Bernstein 2005, zgłoszona do eSTREAM (europejski konkurs na szyfry strumieniowe, 2004 do 2008). W końcowym portfelu z 2008 roku jest Salsa20/12, czyli wersja z 12 rundami, w profilu ,,software''. Te same klocki: tabela 4 na 4, ARX, 20 rund, dodanie tabeli startowej. Różnice: stałe na przekątnej, nonce 8 bajtów i licznik 8 bajtów (64 bity). ChaCha20 z RFC 8439 ma nonce 12 bajtów i licznik 4 bajty; oryginalna ChaCha Bernsteina (i libsodium crypto\_stream\_chacha20) miała, jak Salsa, 8 i 8. IETF wydłużył nonce do 12 bajtów dla zgodności z zaleceniami RFC 5116 (RFC 8439, punkt 2.3; np. 4 bajty nadawcy i 8 bajtów licznika), kosztem krótszego licznika (256 GiB na wiadomość).}
\end{frame}

\samedot
\begin{frame}{Two quarter rounds side by side}
  \vfill
  \begin{columns}[T,onlytextwidth]
    \column{0.48\textwidth}
      \centering\textbf{Salsa20}\\[4pt]
      \salsaqrcode\\[6pt]
      \small each word changes \textbf{once}\\ (b, then c, d, a)\\[3pt]
      rounds: columns, then rows
    \column{0.48\textwidth}
      \centering\textbf{ChaCha20}\\[4pt]
      \qrcode{0}\\[6pt]
      \small each word changes \textbf{twice}\\ (a, d, c, b, and again)\\[3pt]
      rounds: columns, then diagonals
  \end{columns}
  \vfill
  \begin{center}\small Both: 12 small steps on 4 words. Which one spreads a change further?\end{center}
  \vfill
  \keymsg{The same amount of work per quarter round, but ChaCha20 touches every word twice.}
  \say{Ćwiartka rundy obu szyfrów obok siebie. Tyle samo pracy: 4 dodawania, 4 XOR-y, 4 rotacje. W Salsie każde słowo zmienia się raz, w ChaCha dwa razy, więc zmiana jednego bitu powinna rozejść się szybciej. Rotacje 16 i 8 w ChaCha to przesunięcia o całe bajty, szybkie na wielu procesorach. Układ ChaCha (kolumny, potem przekątne, stałe w pierwszym wierszu) wygodnie liczy się instrukcjami wektorowymi: cały wiersz to jeden wektor. Pytanie na dole sprawdzamy w zadaniu 3.}
\end{frame}

\begin{frame}{Task 3: which one mixes faster?}
  \vfill
  \taskgoal{flip one bit and count how far it spreads in Salsa20 and in ChaCha20.}
  \taskready{\texttt{salsa.py} and \texttt{chacha.py}: \texttt{quarter\_round}, \texttt{rounds}, \texttt{initial\_state}.}
  \vfill
  \opentask{Task3}{about 8 minutes}
  \vfill
  \say{Około 8 minut. Najpierw przewidywanie: która ćwiartka rozniesie zmianę dalej? Studenci odwracają jeden bit wejścia, liczą zmienione bity wyjścia (popcount z XOR), potem średnio po 128 bitach, potem cała tabela po 1 do 4 rundach. Moduły są gotowe; kroki w README.}
\end{frame}

\samedot
\begin{frame}{Task 3: what we saw}
  % CHECKED: Task3/solution/mix.py, python3 3.12: one quarter round, bit 0 of a: Salsa20 31, ChaCha20 36;
  %          average over 128 bits 15.8 and 25.1; whole table, key bit 0 flipped, of 512 bits:
  %          1 round 15 / 37, 2 rounds 124 / 245, 3 rounds 234 / 238, 4 rounds 265 / 256
  \vfill
  \begin{center}
    \small one quarter round, average bits changed (of 128):\quad Salsa20 \textbf{15.8}\quad ChaCha20 \textbf{25.1}\\[8pt]
    \begin{tikzpicture}
      \node[font=\scriptsize\bfseries, text=nInk, anchor=south] at (2.6,2.95) {whole table: bits changed (of 512) after $n$ rounds};
      \foreach \n/\a/\b in {1/15/37, 2/124/245, 3/234/238, 4/265/256} {
        \fill[nLine] ({1.3*\n-1.0},0) rectangle ({1.3*\n-0.55},{\a/100});
        \fill[AccentOrange] ({1.3*\n-0.5},0) rectangle ({1.3*\n-0.05},{\b/100});
        \ifnum\a<100 \def\anc{south}\def\col{nInk}\else\def\anc{north}\def\col{white}\fi
        \node[font=\tiny\bfseries, text=nInk, anchor=\anc, inner sep=2pt] at ({1.3*\n-0.775},{\a/100}) {\a};
        \ifnum\b<100 \def\anc{south}\def\col{nInk}\else\def\anc{north}\def\col{white}\fi
        \node[font=\tiny\bfseries, text=\col, anchor=\anc, inner sep=2pt] at ({1.3*\n-0.275},{\b/100}) {\b};
        \node[font=\scriptsize, text=nInk, anchor=north] at ({1.3*\n-0.525},0) {\n};
      }
      \draw[nLine] (0,0) -- (5.4,0);
      \draw[DeepBlue, thick, dashed] (0,2.56) -- (5.4,2.56) node[right, font=\scriptsize, text=DeepBlue] {random: 256};
      \node[font=\scriptsize, text=nInk, anchor=north] at (2.6,-0.35) {rounds};
      \fill[nLine] (7.6,1.6) rectangle (7.9,1.85); \node[font=\scriptsize, anchor=west] at (8.0,1.72) {Salsa20};
      \fill[AccentOrange] (7.6,1.1) rectangle (7.9,1.35); \node[font=\scriptsize, anchor=west] at (8.0,1.22) {ChaCha20};
    \end{tikzpicture}
  \end{center}
  \vfill
  \keymsg{ChaCha20 mixes more per round: random-looking after 2 rounds, Salsa20 after 3 to 4.}
  \say{Wyniki z rozwiązania. Jeden bit to jedna próbka (31 i 36, mała różnica), dlatego liczymy średnią: 15,8 kontra 25,1 bitu na ćwiartkę. W całej tabeli ChaCha20 po 2 rundach jest już przy połowie bitów (245 z 512, czyli jak losowo), Salsa20 dopiero po 3 lub 4. Oba robią 20 rund, więc oba mają duży zapas; ChaCha20 dostaje ten sam zapas szybciej.}
\end{frame}

% CHECKED: libsodium crypto_secretbox = XSalsa20 + Poly1305 (doc.libsodium.org/secret-key_cryptography/secretbox, 2026-11-19);
%          scrypt uses the Salsa20/8 core (tested: salsa_check.py against hashlib.scrypt); secondary sources only: best published attacks: Salsa20/8 (Aumasson et al. 2008,
%          about 2^251; later about 2^245) and ChaCha 7 rounds (about 2^248, later about 2^243): eprint.iacr.org/2015/698,
%          aumasson.jp/data/papers/AFKMR08.pdf
\begin{frame}{Why ChaCha20 won, and where Salsa20 lives on}
  \vfill
  \begin{columns}[T,onlytextwidth]
    \column{0.5\textwidth}
      \textbf{ChaCha20 is the better version}\\[6pt]
      \small
      \ok\ more mixing per round (Task 3)\\[4pt]
      \ok\ a row of the table is one vector:\\ \hspace*{1.1em}fast with SIMD instructions\\[4pt]
      \ok\ best attacks: 7 of 20 rounds\\ \hspace*{1.1em}(Salsa20: 8 of 20); both far from practical
    \column{0.46\textwidth}
      \textbf{Salsa20 is still in use}\\[6pt]
      \small
      \textcolor{nInk}{\faBox}\ \ NaCl and libsodium: \texttt{crypto\_secretbox}\\ \hspace*{1.4em}= XSalsa20 + Poly1305 (the tag of Section 5)\\[4pt]
      \textcolor{nInk}{\faKey}\ \ scrypt, a password hash (Lecture 7):\\ \hspace*{1.4em}built on Salsa20 with 8 rounds
  \end{columns}
  \vfill
  \curious{chacha20/\#6}
  \keymsg{New designs pick ChaCha20; Salsa20 lives on where it was chosen first.}
  \say{Dlaczego ChaCha20 wygrała: więcej mieszania na rundę (zadanie 3), wygodniejsza dla instrukcji wektorowych (SIMD: to samo działanie na kilku słowach naraz), a w kryptoanalizie zapas jest podobny albo lepszy: najlepsze opublikowane ataki sięgają 8 rund Salsy (Aumasson i in. 2008) i 7 rund ChaCha, ze złożonością rzędu 2 do potęgi 240 i więcej, czyli tylko trochę poniżej przeszukania wszystkich kluczy. (Preprint z 2025 roku obiecuje dużo mniej, ale nie był recenzowany.) Salsa20 żyje dalej: w bibliotekach NaCl i libsodium funkcja crypto\_secretbox to XSalsa20-Poly1305 (XSalsa20 to Salsa z 24-bajtowym nonce, Bernstein 2008, wzór dla XChaCha20). scrypt (wykład 7) używa rdzenia Salsa20/8 do mieszania pamięci.}
\end{frame}

\begin{frame}{See it: does the key schedule remember the key?}
  \opensim{chacha20/\#5}
  \vfill
  \begin{enumerate}\small
    \item Card 1, 100\,000 random keys: is byte 0 of what the key setup gives equal to \texttt{key[0]}? Press \textbf{run}.
    \item RC4's table S: far too often. ChaCha20's block: as often as chance. Tick \textbf{seed 2026} for these numbers.
  \end{enumerate}
  \vfill
  \begin{center}
    % CHECKED: random.seed(2026), 100 000 keys: RC4 36.75%, ChaCha20 0.40% (chance 0.39%)
    \begin{tikzpicture}
      \fill[AccentOrange] (0,0.7) rectangle ({36.75*0.14},1.15);
      \node[font=\small\bfseries, text=nInk, anchor=east, align=right] at (-0.15,0.92) {RC4: S after the key schedule\\[-2pt] {\mdseries\scriptsize\texttt{S[0] == key[0]}}};
      \node[font=\small, text=nInk, anchor=west] at ({36.75*0.14+0.15},0.92) {\textbf{36.75\%}\ \ $\leftarrow$ remembers};
      \fill[nLine] (0,-0.22) rectangle ({0.40*0.14},0.22);
      \node[font=\small\bfseries, text=nInk, anchor=east, align=right] at (-0.15,0) {ChaCha20: the output block\\[-2pt] {\mdseries\scriptsize\texttt{byte 0 == key[0]}}};
      \node[font=\small, text=nInk, anchor=west] at (0.25,0) {\textbf{0.40\%}\ \ $\leftarrow$ like chance (1/256 = 0.39\%)};
      \draw[nLine] (0,-0.45) -- (0,1.4);
    \end{tikzpicture}
  \end{center}
  \vfill
  \fixtag{2}%
  \keymsg{That was RC4 problem 2 (a). ChaCha20's block shows nothing of the key.}
  \say{Pokazać na żywo w symulatorze, karta pierwsza. Uczciwie: u RC4 mierzymy tablicę S po układaniu klucza (KSA), a nie bajt wyjścia; u ChaCha20 mierzymy gotowy blok wyjścia, czyli to, co z ChaCha20 wychodzi. To, co RC4 przygotowuje z klucza, w 37\% przypadków ma S[0] = key[0] (jak w Task 7 z wykładu 1). W WEP ten stan przecieka do pierwszych bajtów keystreamu i na tym stoi atak FMS. Blok ChaCha20 trafia w key[0] tylko w 0,4\% kluczy, czyli jak przypadek. To był problem 2, objaw (a). Z zaznaczonym ziarnem 2026 wychodzą liczby ze slajdu.}
\end{frame}

\samedot
\begin{frame}{See it: the second byte}
  \opensim{chacha20/\#5}
  \vfill
  \begin{enumerate}\small
    \item Card 2: 100\,000 random keys: how often is the second keystream byte 0? Press \textbf{run}.
    \item RC4 hits 0 twice too often; ChaCha20 stays at chance. Tick \textbf{seed 2026} for the numbers below.
  \end{enumerate}
  \vfill
  \begin{center}
    % CHECKED: random.seed(2026), 100 000 keys: RC4 0.78% (779 keys), ChaCha20 0.36% (362 keys);
    %          362 is within the noise of 100 000 / 256 = 391 (one standard deviation is about 20)
    \begin{tikzpicture}
      \fill[AccentOrange] (0,0.7) rectangle ({0.78*10},1.15);
      \node[font=\small\bfseries, text=nInk, anchor=east] at (-0.15,0.92) {RC4};
      \node[font=\small, text=nInk, anchor=west] at ({0.78*10+0.15},0.92) {\textbf{0.78\%}\ \ (twice too often)};
      \fill[nLine] (0,-0.22) rectangle ({0.36*10},0.22);
      \node[font=\small\bfseries, text=nInk, anchor=east] at (-0.15,0) {ChaCha20};
      \node[font=\small, text=nInk, anchor=west] at ({0.39*10+0.25},0) {\textbf{0.36\%}\ \ (like chance: \textcolor{DeepBlue}{1/256 = 0.39\%})};
      \draw[DeepBlue, very thick, dashed] (3.906,-0.4) -- (3.906,1.3);
      \draw[nLine] (0,-0.45) -- (0,1.4);
    \end{tikzpicture}
  \end{center}
  \vfill
  \fixtag{2}%
  \keymsg{RC4 leaked two ways; ChaCha20's mixing closes both.}
  \say{Pokazać na żywo w symulatorze, karta druga. To był problem 2, objaw (b): drugi bajt bywa 0 dwa razy za często. ChaCha20 daje 0,36\%: to szum losowania, przy 100\,000 kluczy spodziewamy się około 390 zer, plus minus 20. Domknięcie: oba objawy to jeden mechanizm, dokładne mieszanie, zmierzony raz od strony klucza, raz od strony wyjścia.}
\end{frame}

\begin{frame}{How do we know it is secure?}
  \vfill
  \begin{center}\small Our tests can only \textbf{catch} a bad cipher (RC4). They cannot \textbf{prove} a good one.\\[3pt]
    Real confidence: experts attack it for years and get stuck.\end{center}
  \vfill
  \begin{center}
    \begin{tikzpicture}
      \node[font=\small\bfseries, text=nInk, anchor=east] at (-0.2,0) {rounds:};
      \foreach \k in {1,...,20} {
        \ifnum\k<8
          \node[draw=aWrite, fill=aWrite!25, minimum width=0.5cm, minimum height=0.55cm, font=\tiny, text=nInk, inner sep=0pt] at ({0.56*(\k-1)},0) {\k};
        \else
          \node[draw=aExec, fill=aExec!12, minimum width=0.5cm, minimum height=0.55cm, font=\tiny, text=nInk, inner sep=0pt] at ({0.56*(\k-1)},0) {\k};
        \fi
      }
      \draw[decorate, decoration={brace, mirror, amplitude=5pt}, aWrite, thick] (-0.25,-0.38) -- ({0.56*6+0.25},-0.38)
        node[midway, below=6pt, font=\scriptsize\bfseries, text=aWrite, align=center] {best known attacks reach 7,\\ far from practical};
      \draw[decorate, decoration={brace, mirror, amplitude=5pt}, aExec, thick] ({0.56*7-0.25},-0.38) -- ({0.56*19+0.25},-0.38)
        node[midway, below=6pt, font=\scriptsize\bfseries, text=aExec!70!black] {13 of margin};
      \draw[->, very thick, DeepBlue] ({0.56*19},1.05) -- ({0.56*19},0.35);
      \node[font=\scriptsize\bfseries, text=DeepBlue, anchor=south] at ({0.56*19},1.05) {ChaCha20 runs 20};
    \end{tikzpicture}
  \end{center}
  \vfill
  \fixtag{2}%
  \keymsg{We trust ChaCha20 because attacks stall far short of 20 rounds.}
  \say{Nasze testy potrafią tylko wykryć wadę (jak w RC4), nie zaświadczyć, że dobry szyfr jest bezpieczny. Pewność bierze się z kryptoanalizy: najlepsze znane ataki sięgają 7 z 20 rund, a i tak są tylko teoretyczne (rzędu $2^{243}$ do $2^{248}$ operacji, niewiele mniej niż przeszukanie wszystkich $2^{256}$ kluczy), więc zostaje duży zapas. Dla Salsy20 to 8 rund. ChaCha8 i ChaCha12 to mniej rund: szybsze, mniejszy zapas. Bez szczegółów ataków różnicowych i liniowych.}
\end{frame}

% ======================================================================
\section[Fast without special hardware]{Fast without special hardware}
% ======================================================================

\begin{frame}{AES-CTR: the same trick}
  \begin{itemize}
    \item<1-> \textbf{AES} is a \textbf{block cipher}: it scrambles a fixed 16-byte block.
  \end{itemize}
  \vspace{0.1cm}
  \begin{center}
    \begin{tikzpicture}[>=Stealth]
      \foreach \k in {0,...,15} \node[ptb, minimum width=0.22cm, minimum height=0.4cm, inner sep=0pt] at ({0.24*\k},0) {};
      \node[font=\scriptsize, text=nInk, anchor=west] at (-0.6,0) {in};
      \node[hwmain, minimum width=1.4cm, minimum height=0.7cm, font=\small\bfseries] (aes) at (5.3,0) {AES};
      \node[actorlight=sAmber, font=\tiny\bfseries, inner sep=2pt] (ak) at (5.3,0.75) {\faKey\ key};
      \draw[data] (ak) -- (aes);
      \draw[data] (3.85,0) -- (aes.west); \draw[data] (aes.east) -- (6.75,0);
      \foreach \k in {0,...,15} \node[ctb, minimum width=0.22cm, minimum height=0.4cm, inner sep=0pt] at ({6.95+0.24*\k},0) {};
      \node[font=\scriptsize, text=nInk, anchor=west] at (10.85,0) {out};
      \node[font=\scriptsize, text=gray, anchor=south] at (1.8,0.25) {16 bytes};
      \node[font=\scriptsize, text=gray, anchor=south] at (8.75,0.25) {16 bytes};
    \end{tikzpicture}
  \end{center}
  \vfill
  \begin{itemize}
    \item<2-> In \textbf{CTR mode} a counter turns AES into a keystream: the same trick\\
      as the counter we used for ChaCha20 (Section 2).
  \end{itemize}
  \uncover<2->{%
  \begin{center}
    \begin{tikzpicture}[>=Stealth]
      \foreach \c in {0,1,2} {
        \node[draw=sAmber, line width=1pt, dashed, fill=white, rounded corners=2pt, font=\scriptsize\bfseries, inner sep=2pt] (t\c) at ({3.0*\c},0) {counter \c};
        \node[hwmain, minimum width=1.1cm, minimum height=0.5cm, font=\scriptsize\bfseries] (a\c) at ({3.0*\c+1.55},0) {AES};
        \node[ksb, minimum width=2.6cm, minimum height=0.42cm, font=\scriptsize] (z\c) at ({3.0*\c+1.0},-0.75) {16 bytes};
        \draw[data] (t\c) -- (a\c); \draw[data] (a\c.south) -- (a\c.south |- z\c.north);
      }
      \node[font=\large, text=nInk] at (8.9,-0.75) {\ldots};
      \node[font=\scriptsize, text=nInk, anchor=west] at (9.3,-0.75) {keystream};
    \end{tikzpicture}
  \end{center}}
  \vfill
  \uncover<2->{\keymsg{AES-CTR and ChaCha20 are both counter-based keystreams.}}
  \say{AES sam w sobie szyfruje blok 16 bajtów (szyfr blokowy). W trybie CTR licznik robi z niego keystream, dokładnie tak jak nasz licznik w ChaCha20 z sekcji 2. Więc AES-CTR i ChaCha20 to ten sam rodzaj narzędzia.}
\end{frame}

\samedot
\begin{frame}{ChaCha20 or AES-CTR?}
  \vfill
  \uncover<1->{\begin{center}\large Same job, two engines: which one?\end{center}}
  \vfill
  \begin{columns}[c,onlytextwidth]
    \column{0.5\textwidth}
      \uncover<2->{%
      \centering
      \begin{tikzpicture}\chip{0,0}{1}\end{tikzpicture}\\[4pt]
      \small \textbf{AES-CTR}: fast and safe\\ \textbf{when} the CPU has an AES instruction (AES-NI)}
    \column{0.5\textwidth}
      \uncover<3->{%
      \centering
      \begin{tikzpicture}\chip{0,0}{0}\end{tikzpicture}\\[4pt]
      \small \textbf{ChaCha20}: fast and safe\\ on \textbf{any} CPU}
  \end{columns}
  \vfill
  \uncover<4->{\keymsg{The real contest is ChaCha20 vs AES-CTR, and it depends on the CPU.}}
  \say{To samo zadanie, dwa silniki. AES-CTR wygrywa, gdy procesor ma instrukcje AES. ChaCha20 działa wszędzie. Reszta sekcji pokazuje, co się dzieje bez tych instrukcji.}
\end{frame}

\begin{frame}{Software AES can leak the key through timing}
  \begin{columns}[c,onlytextwidth]
    \column{0.5\textwidth}
      \centering
      \begin{tikzpicture}[>=Stealth]
        \node[font=\scriptsize\bfseries, text=nInk] at (0.3,1.55) {AES in software};
        \foreach \k in {0,...,7} \node[cell, minimum width=1.5cm, minimum height=0.3cm] (t\k) at (0.3,{1.0-0.32*\k}) {};
        \node[font=\tiny, text=gray, anchor=north] at (0.3,{1.0-0.32*7-0.15}) {a table in memory};
        \draw[cellmark] (t5.south west) rectangle (t5.north east);
        \node[actorlight=sAmber, font=\scriptsize\bfseries] (k) at (-1.9,{1.0-0.32*5}) {\faKey\ key};
        \draw[data] (k) -- (t5.west) node[midway, above, font=\tiny, text=nInk] {which entry?};
        \node[text=AccentOrange, font=\huge] at (2.0,0.1) {\faStopwatch};
        \node[font=\tiny, text=AccentOrange, anchor=north] at (2.0,-0.35) {fast or slow?};
      \end{tikzpicture}
    \column{0.48\textwidth}
      \centering
      \begin{tikzpicture}[>=Stealth]
        \node[font=\scriptsize\bfseries, text=nInk] at (0,1.55) {ChaCha20};
        \node[hwmain, minimum width=2.6cm, minimum height=1.3cm, font=\large\bfseries] at (0,0.45) {+\ \ <\kern-1pt<\kern-1pt<\ \ $\oplus$};
        \node[font=\scriptsize, text=nInk] at (0,-0.55) {no table to read};
        \node[font=\scriptsize, text=aExec!70!black] at (0,-0.9) {\ok\ the same steps for every key};
      \end{tikzpicture}
  \end{columns}
  \vfill
  \begin{center}\small Without an AES chip, fast AES uses \textbf{lookup tables}. Which entry it reads depends on the key.\\[2pt]
    An attacker who measures the timing can work the key back out.\\[4pt]
    ChaCha20 (add, rotate, xor): no tables. The same steps every time, whatever the key: nothing to leak.\\[4pt]
    It needs only plain integer instructions, and it fits vector (SIMD) instructions on any CPU.\end{center}
  \vfill
  \keymsg{A key-dependent table lookup can leak the key; ARX has none.}
  \say{Szybki AES w programie czyta tablice, a to, którą komórkę, zależy od klucza. Napastnik mierzy czas dostępu (cache) i odtwarza klucz. To jest atak czasowy. ChaCha20 robi dodawanie, rotację i XOR zawsze tak samo, niezależnie od klucza, więc nie ma czego podsłuchać. (Da się zrobić stałoczasowy AES, ale to osobny, sprytny kod SIMD.) SIMD to instrukcje wektorowe: to samo działanie na kilku słowach naraz; ChaCha20 dobrze się na nie układa, bo cały wiersz tabeli to jeden wektor. Źródła ataku: Bernstein 2005 (,,Cache-timing attacks on AES''), Osvik, Shamir, Tromer 2005.}
\end{frame}


\begin{frame}{Measured on a laptop}
  \begin{center}\small\textcolor{gray}{Intel Core i5-8265U\ \ $\cdot$\ \ OpenSSL 3.0.13\ \ $\cdot$\ \ 16 KiB blocks}\end{center}
  \vfill
  \begin{center}
    % CHECKED: 2026-11-25, the author's laptop (i5-8265U, OpenSSL 3.0.13), openssl speed -evp X -bytes 16384 -seconds 2,
    %   two runs (local/audyt/tmp-L02/speed/out2.txt): AES-128-CTR 6.03 / 5.71 GB/s, ChaCha20 3.22 / 2.99 GB/s;
    %   OPENSSL_ia32cap="~0x200000200000000:~0x0" (AES-NI and PCLMULQDQ off): AES-128-CTR 0.62 / 0.59, ChaCha20 3.14 / 3.12.
    %   Without ":~0x0" OpenSSL also forgets AVX2 and ChaCha20 drops to 1.6: not the question we ask.
    \begin{tikzpicture}
      \node[font=\small\bfseries, text=nInk] at (2.3,2.2) {with AES-NI};
      \node[font=\small\bfseries, text=nInk] at (8.3,2.2) {without AES-NI};
      \foreach \x/\a/\c in {0/5.9/3.1, 6.0/0.6/3.1} {
        \fill[nLine] (\x,1.1) rectangle ({\x+0.7*\a},1.6);
        \node[font=\small, text=nInk, anchor=west] at ({\x+0.7*\a+0.1},1.35) {\a\ GB/s};
        \fill[AccentOrange] (\x,0.3) rectangle ({\x+0.7*\c},0.8);
        \node[font=\small, text=nInk, anchor=west] at ({\x+0.7*\c+0.1},0.55) {\c\ GB/s};
        \draw[nLine] (\x,0.1) -- (\x,1.85);
      }
      \node[font=\small\bfseries, text=nInk, anchor=east] at (-0.15,1.35) {AES-128-CTR};
      \node[font=\small\bfseries, text=nInk, anchor=east] at (-0.15,0.55) {ChaCha20};
      \draw[annot] (5.85,0.2) rectangle (9.85,0.9);
    \end{tikzpicture}
  \end{center}
  \vfill
  \begin{center}\scriptsize\textcolor{gray}{without AES-NI = the same laptop with the AES instructions switched off; both measured}\end{center}
  \keymsg{With AES chips, AES wins. Without them, ChaCha20 wins.}
  \say{Prawdziwe liczby z mojego laptopa (i5-8265U, openssl speed, pomiar 25.11.2026, średnia z dwóch uruchomień). Ze sprzętem AES wygrywa około 2 razy. Z wyłączonymi instrukcjami AES (zmienna OPENSSL\_ia32cap) AES spada do 0,6 GB/s, a ChaCha20 się nie zmienia, bo tych instrukcji nie używa: wygrywa około 5 razy. Oba słupki ChaCha20 są zmierzone. (Uwaga, gdyby ktoś powtarzał pomiar: maska bez drugiej części ,,:\textasciitilde 0x0'' gasi w OpenSSL także AVX2 i wtedy ChaCha20 spada do 1,6 GB/s.) Zegar laptopa się zmienia, więc liczby też, proporcje zostają.}
\end{frame}

\begin{frame}{Why Google added ChaCha20 to HTTPS}
  \begin{columns}[c,onlytextwidth]
    \column{0.24\textwidth}
      \centering
      \begin{tikzpicture}
        \draw[nInk, very thick, fill=white, rounded corners=6pt] (0,0) rectangle (1.8,3.2);
        \draw[nLine, fill=nFill] (0.15,2.55) rectangle (1.65,2.9);
        \node[font=\tiny, text=nInk] at (0.9,2.72) {\faLock\ https://};
        \foreach \y in {2.2,1.9,1.6,1.3} \draw[nLight, thick] (0.25,\y) -- (1.55,\y);
        \node[draw=white, dashed, fill=nInk!70, rounded corners=1pt, font=\tiny, text=white, minimum width=1.2cm] at (0.9,0.75) {no AES};
        \fill[nInk] (0.9,0.25) circle (0.1);
      \end{tikzpicture}
    \column{0.74\textwidth}
      \small
      \textbf{2014}: many Android phones had no hardware AES.\\[3pt]
      HTTPS with AES was slow, so Google added \textbf{ChaCha20-Poly1305} to Chrome.\\[8pt]
      \textbf{Today}: most phones have hardware AES and often use AES-GCM (AES with a tag).\\[3pt]
      But ChaCha20-Poly1305 stayed a standard HTTPS cipher (TLS 1.3),\\ chosen wherever there is no AES hardware.\\[8pt]
      ``Google uses it'' = browser and server pick it per HTTPS connection, not the whole phone.
  \end{columns}
  \vfill
  \keymsg{No AES hardware $\rightarrow$ ChaCha20; it stayed a standard either way.}
  % CHECKED: Google Security Blog, April 2014 (local/audyt/fakty.md); TLS 1.3: RFC 8446 section 9.1
  \say{W 2014 wiele telefonów z Androidem nie miało sprzętowego AES, więc HTTPS przez AES był wolny; Google dodał ChaCha20-Poly1305 do Chrome (BoringSSL) i swoich serwerów. Chrome wybiera adaptacyjnie: gdy procesor ma sprzętowy AES, woli AES-GCM, a bez niego ChaCha20. Systemowa biblioteka TLS Androida (Conscrypt) też to obsługuje, więc nie tylko Chrome. ,,Google używa'' znaczy: przeglądarka i serwer negocjują szyfr dla danego połączenia HTTPS, a nie ,,cały telefon''. Osobny przypadek: szyfrowanie dysku tanich telefonów bez AES to Adiantum (XChaCha12), inny użytek niż HTTPS. (Źródła: Google Security Blog, kwiecień 2014; RFC 8446, sekcja 9.1: TLS\_CHACHA20\_POLY1305\_SHA256 jako SHOULD.)}
\end{frame}

% ======================================================================
\section[ChaCha20-Poly1305]{ChaCha20-Poly1305}
% ======================================================================

\begin{frame}{Recall: encryption does not protect integrity}
  \vfill
  \begin{center}
    % CHECKED: python3, aead.chacha20 with the key and nonce of Task 4, counter 1: "Pay Bob 100" gives the first row;
    %          byte 8 XOR 0x08 (D7 -> DF) decrypts to "Pay Bob 900"
    \tamperfig{4C,EA,76,9D,37,9F,55,B9,D7,A3,5A}{4C,EA,76,9D,37,9F,55,B9,DF,A3,5A}
  \end{center}
  \vfill
  \begin{center}\small Bob decrypts and reads a changed message, and cannot tell.\\[3pt]
    Same problem as RC4: a stream cipher hides, it does not protect.\end{center}
  \vfill
  \fixtag{3}%
  \keymsg{We need a way to detect any change.}
  \say{Przypomnienie z sekcji 1: goły szyfr nie chroni przed zmianą. Tu już ChaCha20 zamiast RC4, a efekt ten sam: 100 zmienia się w 900. To ostatnia naprawa.}
\end{frame}

\samedot
\begin{frame}{A tag}
  \begin{center}\small \textbf{Poly1305} computes a 16-byte \textbf{tag} from the ciphertext and a secret key.\end{center}
  \vfill
  \begin{center}
    % CHECKED: Task4/aead.py, seal(key, nonce, b"Pay Bob 100") with the key and nonce of Task 4:
    %          ciphertext 4c ea 76 9d 37 9f 55 b9 d7 a3 5a, tag 0e e9 c8 47 ... 30 71;
    %          with byte 8 XOR 08 ("Pay Bob 900", the slide before) the recomputed tag is 0c 11 a0 78 ... 04 3e
    \begin{tikzpicture}[>=Stealth]
      \node[font=\small\bfseries, text=nInk, anchor=east] at (-0.4,0) {send:};
      \byterow{0}{0}{ctb}{4C,EA,76,9D,37,9F,55,B9,D7,A3,5A}{0.66}
      \node[draw=nInk, very thick, fill=nFill, rounded corners=1.5pt, minimum width=3.2cm, minimum height=0.48cm, font=\ttfamily\scriptsize, text=nInk] (tag) at (8.95,0) {0e e9 c8 47 \ldots\ 30 71};
      \node[font=\scriptsize, text=nInk, anchor=south] at (3.3,0.3) {ciphertext (11 bytes)};
      \node[font=\scriptsize, text=nInk, anchor=south] at (8.95,0.3) {\faStamp\ tag (16 bytes)};
      \node[font=\small\bfseries, text=nInk, anchor=east] at (-0.4,-1.0) {Bob:};
      \node[font=\small, text=nInk, anchor=west] at (-0.2,-1.0) {recomputes the tag from what arrived};
      \node[font=\small, text=nInk, anchor=west] at (-0.2,-1.5) {\ok\ unchanged: \texttt{0e e9 c8 47 \ldots} = the tag: accept};
      \node[font=\small, text=nInk, anchor=west] at (-0.2,-1.95) {\nok\ Mallory's \texttt{900}: \texttt{0c 11 a0 78 \ldots} $\neq$ the tag: \textbf{reject}};
    \end{tikzpicture}
  \end{center}
  \vfill
  \bigdef{\textbf{Authenticated encryption}: hide the message \textbf{and} detect any change.}
  \fixtag{3}%
  \keymsg{A wrong tag means a changed message: throw it away.}
  \say{Poly1305 dolicza 16-bajtowy znacznik. Odbiorca liczy go ponownie; jeśli się różni, odrzuca wiadomość. Bez klucza Mallory nie umie policzyć znacznika dla zmienionej wiadomości (dlaczego, na następnym slajdzie). Bajty na slajdzie to prawdziwy wynik aead.py z zadania 4, dla zmiany 100 na 900 z poprzedniego slajdu; zadanie sprawdza inną zmianę.}
\end{frame}

\samedot
\begin{frame}{How Poly1305 makes the tag}
  \vfill
  \begin{center}
    \begin{tikzpicture}[>=Stealth]
      \foreach \k/\l in {0/1, 1/2, 2/3} {
        \node[ctb, minimum width=2.3cm, minimum height=0.55cm, font=\scriptsize] (m\k) at ({3.1*\k},1.4) {chunk $m_\l$ (16 bytes)};
      }
      \node[font=\scriptsize, text=nInk, anchor=west] at (7.6,1.4) {the ciphertext, cut into 16-byte chunks};
      \node[draw=nInk, thick, fill=nFill, rounded corners=2pt, font=\scriptsize\ttfamily, minimum height=0.55cm] (a0) at (-1.7,0) {acc = 0};
      \foreach \k/\l in {0/1, 1/2, 2/3} {
        \node[hwmain, minimum width=2.3cm, minimum height=0.55cm, font=\scriptsize] (s\k) at ({3.1*\k},0) {(acc + $m_\l$) $\cdot$ \textbf{r}};
        \draw[data] (m\k) -- (s\k);
      }
      \draw[data] (a0) -- (s0); \draw[data] (s0) -- (s1); \draw[data] (s1) -- (s2);
      \node[circle, draw=AccentOrange, very thick, fill=white, font=\small\bfseries, text=AccentOrange, inner sep=1pt] (pl) at (8.4,0) {+ \textbf{s}};
      \draw[data] (s2) -- (pl);
      \node[draw=nInk, very thick, fill=nFill, rounded corners=1.5pt, font=\scriptsize\bfseries, minimum height=0.55cm] (tg) at (10.6,0) {\faStamp\ tag (16 bytes)};
      \draw[data] (pl) -- (tg);
      \node[font=\scriptsize, text=nInk, anchor=north] at (3.1,-0.45) {every step modulo the prime $p = 2^{130} - 5$};
    \end{tikzpicture}
  \end{center}
  \vfill
  \begin{center}\small $\text{tag} = (m_1 r^3 + m_2 r^2 + m_3 r \bmod p) + s$:\ a polynomial in the secret point $r$, then hidden by $s$.\\[3pt]
    Mallory changes a chunk: the tag changes in a way she cannot compute without $r$ and $s$. She can only guess.\end{center}
  \vfill
  \curious{chacha20/\#7}
  \fixtag{3}%
  \keymsg{The tag mixes every chunk with the secret $r$; without $r$ and $s$, no right tag.}
  \say{Intuicja Poly1305 bez dowodu. Szyfrogram dzielimy na kawałki po 16 bajtów i traktujemy je jak liczby (do każdego dopisuje się bajt 1, żeby kawałek z zerami na końcu różnił się od krótszego). Akumulator: dodaj kawałek, pomnóż przez tajne r, wszystko modulo liczba pierwsza $2^{130}-5$. Wychodzi wielomian od r o współczynnikach z wiadomości. Na koniec dodajemy tajne s. Zmiana dowolnego kawałka zmienia wielomian, a nowej wartości nie da się policzyć bez r; s ukrywa wynik, więc z tagu nie widać r. Szansa trafienia tagu na ślepo jest znikoma (rzędu $2^{-100}$ na próbę dla krótkich wiadomości). W ChaCha20-Poly1305 do kawałków dochodzą jeszcze długości i dane powiązane (zaraz). Symulator, krok 7: akumulator kawałek po kawałku.}
\end{frame}

\samedot
\begin{frame}{Where the tag key comes from}
  \vfill
  \begin{center}
    \begin{tikzpicture}[>=Stealth]
      \node[draw=AccentOrange, thick, fill=AccentOrange!15, minimum width=1.8cm, minimum height=0.7cm, font=\scriptsize] (b0) at (0,0) {64 bytes};
      \foreach \n in {1,2,3} \node[ksb, minimum width=1.8cm, minimum height=0.7cm, font=\scriptsize] at ({2.0*\n},0) {64 bytes};
      \foreach \n in {0,1,2,3} \node[font=\scriptsize\bfseries, text=nInk, anchor=south] at ({2.0*\n},0.4) {counter \n};
      \node[font=\large, text=nInk] at (7.4,0) {\ldots};
      \draw[data] (b0.south) -- (0,-1.0) node[below, font=\scriptsize, text=nInk, align=center] {one-time\\ Poly1305 key};
      \draw[decorate, decoration={brace, mirror, amplitude=5pt}, nInk] (1.1,-0.45) -- (7.0,-0.45)
        node[midway, below=6pt, font=\scriptsize, text=nInk] {the message is encrypted from counter 1 on};
    \end{tikzpicture}
  \end{center}
  \vfill
  \begin{center}\small ChaCha20 block with counter 0 $\rightarrow$ its first 32 bytes are $r$ and $s$; the other 32 are not used.\\[3pt]
    Why new for every message: two tags made with the same $(r, s)$ give Mallory two equations,\\ from which she can compute $r$ and $s$, and then the tag of any message.\\[3pt]
    So the data is encrypted from counter 1 on, and block 0 is a fresh key for each nonce.\end{center}
  \vfill
  \fixtag{3}%
  \keymsg{The tag key is a fresh ChaCha20 block, one per message.}
  \say{Skąd klucz do znacznika: blok ChaCha20 z licznikiem 0, z którego bierzemy tylko pierwsze 32 bajty (16 na r, 16 na s). Po co nowy klucz dla każdej wiadomości: Poly1305 to MAC jednorazowy. Dwa tagi z tym samym (r, s) to dwa równania wielomianowe; odejmując je, s znika i zostaje równanie na r, które da się rozwiązać. Z r i s Mallory liczy tag dowolnej wiadomości. Nowy nonce daje nowy blok 0, więc nowe (r, s). Dane szyfrujemy od licznika 1, stąd licznik 1 w zadaniu 1.}
\end{frame}

\samedot
\begin{frame}{Associated data: read in the clear, still protected}
  \vfill
  \begin{center}
    \begin{tikzpicture}[>=Stealth]
      \node[draw=nInk, thick, fill=white, rounded corners=1.5pt, minimum width=3.2cm, minimum height=0.6cm, font=\scriptsize\ttfamily] (h) at (0,0) {to: Bob, packet 7};
      \node[ctb, minimum width=4.2cm, minimum height=0.6cm, font=\scriptsize] (c) at (3.8,0) {ciphertext};
      \node[draw=nInk, very thick, fill=nFill, rounded corners=1.5pt, minimum width=2.4cm, minimum height=0.6cm, font=\scriptsize] (t) at (7.25,0) {\faStamp\ tag};
      \node[font=\scriptsize, text=nInk, anchor=south] at (h.north) {header: not encrypted};
      \node[font=\scriptsize, text=nInk, anchor=south] at (c.north) {encrypted};
      \draw[decorate, decoration={brace, mirror, amplitude=5pt}, AccentOrange, thick] ([yshift=-0.12cm]h.south west) -- ([yshift=-0.12cm]c.south east)
        node[midway, below=6pt, font=\scriptsize\bfseries, text=AccentOrange] {the tag covers both};
    \end{tikzpicture}
  \end{center}
  \vfill
  \begin{center}\small Routers and servers must read the header (where it goes, which packet), so it is not encrypted.\\[3pt]
    Mallory changes it to \texttt{to: Eve}: the tag does not match, Bob rejects the packet.\end{center}
  \vfill
  \bigdef{\textbf{AEAD}: authenticated encryption with associated data.}
  \fixtag{3}%
  \keymsg{Encrypt what is secret; let the tag protect also what must stay readable.}
  \say{AD w AEAD: dane powiązane, czyli nagłówek, który musi być jawny (adres, numer pakietu, numer rekordu w TLS), ale nie może być zmieniony. Szyfrujemy tylko treść, a znacznik liczymy z nagłówka i szyfrogramu razem (i ich długości). Zmieniony nagłówek przy tym samym szyfrogramie: znacznik się nie zgadza. ChaCha20-Poly1305 to AEAD; parametr aad w aead.py to właśnie te dane. Sprawdzimy to w ostatnim kroku zadania 4.}
\end{frame}

\begin{frame}{Task 4: change one byte}
  \vfill
  \taskgoal{seal a message, change one byte or the header, and see what Bob reads with and without the tag.}
  \taskready{\texttt{aead.py}: \texttt{seal}, \texttt{open\_} and \texttt{chacha20}.}
  \vfill
  \opentask{Task4}{about 7 minutes}
  \vfill
  \fixtag{3}%
  \say{Około 7 minut. aead.py gotowy (seal, open\_, chacha20). Podkreślnik w nazwie, bo open to już funkcja Pythona. Studenci zmieniają bajt 5 (inny niż na slajdzie), odszyfrowują bez znacznika, potem próbują open\_, a na końcu zmieniają tylko nagłówek (aad). Najpierw przewidywanie. Szczegóły w README.}
\end{frame}

\samedot
\begin{frame}{Task 4: what we saw}
  \vfill
  % CHECKED: Task4/solution/tamper.py, python3 3.12: without the tag b'Pay Bnb 100'; with the tag and with the changed
  %          header "open_ raises: bad tag: message rejected"; aead.py gives the ciphertext d31a8d34648e60db... and the
  %          tag 1ae10b594f09e26a7e902ecbd0600691 of RFC 8439 2.8.2
  \begin{center}
    \begin{minipage}{0.72\textwidth}
      \stepbox{Output}{\texttt{without the tag Bob reads: b'Pay Bnb 100'}\\
        \texttt{with the tag: open\_ raises: bad tag: message rejected}\\
        \texttt{header changed: open\_ raises: bad tag: message rejected}}
    \end{minipage}\\[0.4cm]
    Nothing is decrypted: the tag is checked first, in constant time.
  \end{center}
  \vfill
  \fixtag{3}%
  \keymsg{One changed byte: rejected. Without a tag it would pass unseen.}
  \say{Bez znacznika zmiana przechodzi: Bob czyta Pay Bnb 100 i nic nie podejrzewa. Ze znacznikiem open\_ odrzuca pakiet i nic nie odszyfrowuje; tak samo przy zmienionym nagłówku. Dwie zasady implementacji: najpierw sprawdź tag, dopiero potem wydaj tekst jawny; porównuj tagi w stałym czasie (aead.py używa hmac.compare\_digest, bo zwykłe porównanie kończy się na pierwszym różnym bajcie, a czas mógłby zdradzać, ile bajtów podróbki jest dobrych). To jest ostatnia naprawa.}
\end{frame}

\begin{frame}{The nonce rule still holds}
  \vfill
  \begin{center}
    \begin{tikzpicture}[>=Stealth]
      \envelope{e1}{0,0.8}{$C_1$}
      \envelope{e2}{0,-0.8}{$C_2$}
      \node[actorlight=sAmber, font=\scriptsize\bfseries, align=center] at (2.3,0.8) {key, nonce N};
      \node[actorlight=sAmber, font=\scriptsize\bfseries, align=center] at (2.3,-0.8) {key, nonce N};
      \node[font=\scriptsize\bfseries, text=AccentOrange] at (2.3,0) {the same!};
      \draw[data, line width=1.6pt] (3.6,0) -- (4.6,0);
      \node[font=\small, text=nInk, anchor=west, align=left] at (4.7,0) {the same keystream $\rightarrow$ $C_1 \oplus C_2 = \mathbf{P_1 \oplus P_2}$\\[2pt] the same Poly1305 key $\rightarrow$ two tags reveal it: tags can be forged};
    \end{tikzpicture}
  \end{center}
  \vfill
  \begin{center}\small Repeat a nonce under one key: the same keystream.\\[3pt]
    A new cipher does not cancel the old rule.\end{center}
  \vfill
  \fixtag{3}%
  \keymsg{One (key, nonce) pair, one message. Still never twice.}
  \say{Powtórzony nonce to ten sam keystream, czyli dokładnie błąd z sekcji 1, plus ten sam klucz Poly1305 z bloku 0, a dwa tagi z jednym kluczem pozwalają go wyliczyć (poprzedni slajd o kluczu tagu), więc można podrabiać znaczniki. Limity klucza i szyfry odporne na powtórzony nonce (np. AES-GCM-SIV, RFC 8452) to temat na później. ChaCha20 nie zwalnia z reguły ,,never twice''.}
\end{frame}

% CHECKED 2026-11-25, read: RFC 8446 section 5.3 (64-bit record sequence number, padded, XORed with the static write_iv);
%          wireguard.com/protocol (u64 counter, nonce = 32 zero bits + the 64-bit counter); NIST SP 800-38D section 8.3
%          (random IVs: at most 2^32 invocations per key, for AES-GCM);
%          the chance of a repeat after q = 2^32 random 96-bit nonces is about q^2 / 2^97 = 2^-33 (python3: 2**64 / 2**97)
\samedot
\begin{frame}{Choosing nonces in practice}
  \vfill
  \begin{center}
    \renewcommand{\arraystretch}{1.9}
    \small
    \begin{tabular}{c l l}
      \bigbadge{1} & \textbf{count}: nonce = message number & TLS 1.3: a secret IV XOR the record number;\\
                   &                                         & WireGuard: a counter in every packet \\
      \bigbadge{2} & \textbf{one key per direction} & Alice to Bob, Bob to Alice: the counters never meet \\
      \bigbadge{3} & \textbf{random} 12-byte nonce & fine up to about $2^{32}$ messages per key\\
                   &                                & (then a repeat has chance about $2^{-33}$)
    \end{tabular}
  \end{center}
  \vfill
  \fixtag{3}%
  \keymsg{A counter never repeats; a random 12-byte nonce is safe only up to a limit.}
  \say{Jak w praktyce dbać o regułę ,,nigdy dwa razy''. Najprościej licznik: TLS 1.3 bierze tajny IV i XOR-uje go z numerem rekordu, WireGuard wysyła 64-bitowy licznik w każdym pakiecie. Ważne: osobne klucze dla każdego kierunku, bo inaczej licznik Alice i licznik Boba dałyby te same nonce. Losowy nonce 12-bajtowy: z paradoksu urodzin po q wiadomościach szansa powtórki to około $q^2/2^{97}$; NIST dla AES-GCM z losowym IV pozwala na $2^{32}$ wiadomości na klucz, wtedy szansa jest rzędu $2^{-33}$. Kto chce losować bez liczenia, bierze dłuższy nonce: następny slajd.}
\end{frame}

\begin{frame}{Good to know: XChaCha20}
  \vfill
  \begin{center}
    \begin{tikzpicture}
      \node[font=\small\bfseries, text=nInk, anchor=east] at (-0.4,0) {ChaCha20};
      \foreach \k in {0,...,11} \node[draw=sAmber, line width=1pt, fill=white, minimum width=0.42cm, minimum height=0.45cm, inner sep=0pt] at ({0.45*\k},0) {};
      \node[font=\small, text=nInk, anchor=west] at (5.3,0) {12 bytes: random up to $2^{32}$ messages};
      \node[font=\small\bfseries, text=nInk, anchor=east] at (-0.4,-0.9) {XChaCha20};
      \foreach \k in {0,...,23} \node[draw=sAmber, line width=1pt, fill=white, minimum width=0.42cm, minimum height=0.45cm, inner sep=0pt] at ({0.45*\k},-0.9) {};
      \node[font=\small, text=nInk, anchor=west] at (10.7,-0.9) {24 bytes};
      \node[font=\small, text=aExec!70!black, anchor=west] at (5.3,-1.6) {\ok\ random is safe in practice};
    \end{tikzpicture}
  \end{center}
  \vfill
  \fixtag{3}%
  \keymsg{A longer nonce: safe to choose at random.}
  % CHECKED 2026-11-25: datatracker.ietf.org/doc/draft-irtf-cfrg-xchacha: expired, IRTF state Dead, last revision -03 of 2020-01-10
  \say{Wariant XChaCha20 ma 24-bajtowy nonce, więc można go losować bez obaw o powtórkę (nawet po $2^{64}$ wiadomościach szansa jest pomijalna). Jak to działa: HChaCha20 robi z klucza i pierwszych 16 bajtów nonce nowy podklucz, a reszta nonce idzie do zwykłej ChaCha20. Pierwowzorem jest XSalsa20 (Bernstein 2008), ta z libsodium. XChaCha20 nie ma RFC: szkic IRTF (draft-\allowbreak irtf-\allowbreak cfrg-\allowbreak xchacha) wygasł w 2020 roku, ale wariant jest w libsodium i wielu bibliotekach. Slajd opcjonalny.}
\end{frame}

% CHECKED: TLS 1.3, WireGuard, OpenSSH 6.5 (local/audyt/fakty.md); OpenBSD arc4random and Linux 4.8 /dev/urandom use ChaCha20
%          (en.wikipedia.org/wiki/Salsa20, read 2026-11-25; a secondary source)
\begin{frame}{Where you use it}
  \vfill
  \begin{center}
    \begin{tikzpicture}
      \foreach \ic/\t/\s [count=\k from 0] in {\faLock/TLS 1.3/{the padlock in your browser}, \faNetworkWired/WireGuard/{a modern VPN},
          \faTerminal/OpenSSH/{remote login}, \faDice/{Linux, OpenBSD}/{random numbers}} {
        \node[hw, minimum width=2.9cm, minimum height=2.3cm] (b\k) at ({3.3*\k},0) {};
        \node[text=nInk] at ({3.3*\k},0.35) {\fontsize{30}{32}\selectfont\ic};
        \node[font=\normalsize\bfseries, text=nInk] at ({3.3*\k},-0.65) {\t};
        \node[font=\scriptsize, text=gray, below=0.12cm of b\k, text width=2.9cm, align=center] {\s};
      }
    \end{tikzpicture}
  \end{center}
  \vfill
  \fixtag{3}%
  \keymsg{ChaCha20, with Poly1305 or alone, runs on a large part of today's internet.}
  \say{Gdzie dziś żyje: TLS 1.3 (kłódka w przeglądarce), WireGuard (nowoczesny VPN, tylko ChaCha20-Poly1305), OpenSSH (chacha20-poly1305@openssh.com od wersji 6.5, 2014), oraz generatory liczb losowych: OpenBSD zamienił w arc4random RC4 na ChaCha20 , a Linux używa ChaCha20 w /dev/urandom od jądra 4.8 (2016). Ładne domknięcie historii RC4 $\rightarrow$ ChaCha20.}
\end{frame}

% ======================================================================
\section[Side by side]{RC4 and ChaCha20 side by side}
% ======================================================================

\begin{frame}{RC4 and ChaCha20 side by side}
  \vfill
  \begin{center}
    \small
    \renewcommand{\arraystretch}{1.7}
    \begin{tabular}{l l l l}
      \toprule
      \textbf{problem} & \textbf{RC4} & \textbf{ChaCha20-Poly1305} & \textbf{how} \\
      \midrule
      1 nonce and counter inputs? & \nok\ no & \ok\ yes & key + nonce + counter \\
      2 looks random? & \nok\ no (2 ways) & \ok\ yes & 20 rounds of ARX \\
      3 integrity? & \nok\ no & \ok\ yes & Poly1305 tag \\
      fast without an AES chip? & \textcolor{sAmberText}{\faExclamationTriangle}\ yes (broken) & \ok\ yes & ARX \\
      \bottomrule
    \end{tabular}
  \end{center}
  \vfill
  \keymsg{Everything RC4 got wrong, ChaCha20-Poly1305 gets right.}
  \say{Zebranie całego wykładu i spełnienie obietnicy ze slajdu o RC4: trzy problemy (w kolejności napraw), każdy z naprawą, plus szybkość. Każdy wiersz to jedno ,,jak'' po prawej. Pewności co do ChaCha20 nie dały nasze dwa testy (one tylko złapały RC4), tylko lata nieudanych ataków.}
\end{frame}


\end{document}
